% a05_cnum_NEG: proofs of lem:k02:cnum (c-numbers) and prop:k02:NEG ((NEG_infty))
\providecommand{\NEGinf}{\ensuremath{(\mathrm{NEG}_{\infty})}}
\providecommand{\BZ}{\ensuremath{\mathsf B}}
\providecommand{\Bc}{\ensuremath{B^{\mathrm c}}}
\providecommand{\Gh}{\ensuremath{\hat G}}
\providecommand{\intk}{\ensuremath{\int_{k}}}
\providecommand{\intp}{\ensuremath{\int_{p}}}

\section{The $c$-numbers and the growth of the bare coefficient with the mass}\label{app:a05_cnum_NEG}

This appendix proves \cref{lem:k02:cnum} and \cref{prop:k02:NEG}. Both are statements about lattice Feynman integrals on the infinite lattice $\eps\Z^3$: the first about the order-$e^2$ coefficient $X_\eps$ of \cref{def:s02:sources}, which defines the reference scheme $z_k=e^2X_{\eps_k}$, the second about King's free-boundary mass counterterm ${\dm}_\eps(M)$ \cite[(3.32)--(3.33)]{K1}, \cite[p.~330]{K2} at its five orders in $(e^2,\lambda)$. \Cref{app:a05:A} gives the exact asymptotics of $X_\eps$ (\cref{thm:a05:Z}), from which \cref{lem:k02:cnum} follows; \cref{app:a05:B} computes the continuum limit of ${\dm}_\eps(M)-{\dm}_\eps(m)$ in closed form (\cref{thm:a05:E}) and derives from it the lower bound of \cref{prop:k02:NEG} (\cref{thm:a05:B}). Throughout, $N=2$ (two real scalar components, $|\phi|^2=r^2$), the quartic term is $\lambda|\phi|^4$, and the photon is in Feynman form with mass $\lmu$ \cite[(2.4)]{K1}, \cite[(2.5)]{K2}.

\subsection{The order-$e^2$ Wick constant}\label{app:a05:A}

\paragraph{Conventions on the unit lattice.} On $\Z^3$ put $\BZ\defeq[-\pi,\pi]^3$, $\intk\defeq\int_\BZ\frac{d^3k}{(2\pi)^3}$, $\hat q_\alpha^2\defeq4\sin^2(q_\alpha/2)$, $\hat q^2\defeq\sum_\alpha\hat q_\alpha^2$, and for $a,b>0$
\[
\Gh_a(q)\defeq(\hat q^2+a^2)^{-1},\qquad D_b(k)\defeq(\hat k^2+b^2)^{-1},
\]
\begin{gather*}
\hat\Pi_{a,b}(p)\defeq\intk D_b(k)\sum_\alpha\Big[\cos p_\alpha-4\sin^2\tfrac{2p_\alpha+k_\alpha}2\,\Gh_a(p+k)\Big],\\
\bar\pi_{a,b}(p)\defeq\hat\Pi_{a,b}(p)-\hat\Pi_{a,b}(0),\qquad X_1(a,b)\defeq-2\intp\Gh_a(p)^2\,\bar\pi_{a,b}(p).
\end{gather*}

\begin{lemma}[the formula for $X_\eps$]\label{lem:a05:deriv}
\textup{(i)} On $\Z^3$ with King's action at unit spacing, scalar mass $a$, photon mass $b$ and charge $e$, the coefficient of $e^2$ in $\langle|\phi(0)|^2\rangle$ at $\lambda=0$, with the order-$e^2$ mass counterterm fixed by \cite[(3.33)]{K1}, is $X_1(a,b)$. \textup{(ii)} $X_\eps(m;\mu)=X_1(\eps m,\eps\mu)$ for the quantity $X_\eps$ of \cref{def:s02:sources} at scalar mass $m$ and photon mass $\mu$. In particular $z_k=e^2X_1(\eps_km,\eps_k\lmu)$.
\end{lemma}
\begin{proof}
(i) Work on $\Z^3/N\Z^3$ and let $N\to\infty$ at the end (all integrands below are smooth and periodic for $a,b>0$, so the Riemann sums converge). King's matter action at unit spacing is $\frac12\sum_{x,\alpha}|e^{\mathrm ieA_\alpha(x)}\phi(x+e_\alpha)-\phi(x)|^2+\frac12(a^2+{\dm})\sum_x|\phi|^2$ \cite[(2.3), (2.5)]{K2}, i.e.\ $\bar\phi(K+M(A))\phi$ with $K=\frac12(-\Delta+a^2)$ and $M(A)(x,x+e_\alpha)=-\frac12(e^{\mathrm ieA_\alpha(x)}-1)$, $M(A)(x+e_\alpha,x)=-\frac12(e^{-\mathrm ieA_\alpha(x)}-1)$. The photon is Gaussian with $\langle A_\alpha A_\beta\rangle=\delta_{\alpha\beta}(-\Delta+b^2)^{-1}$ (Feynman form). At fixed $A$ the $\phi$-integral gives $\langle\phi(0)\bar\phi(0)\rangle_A=(K+M)^{-1}(0,0)$ and the weight $\det(K+M)^{-1}$. Write $M=eM_1+e^2M_2+O(e^3)$. Then $(K+M)^{-1}=K^{-1}-K^{-1}MK^{-1}+K^{-1}MK^{-1}MK^{-1}-\dots$, and $\Tr K^{-1}M_1=\sum_{x,\alpha}K^{-1}(e_\alpha)(-\frac{\mathrm i}2A_\alpha(x)+\frac{\mathrm i}2A_\alpha(x))=0$ identically, so the determinant does not contribute at order $e^2$. Hence the order-$e^2$ part of $\langle\phi\bar\phi\rangle$ is $K^{-1}\Pi K^{-1}$ with $\Pi=-\langle M_2\rangle+\langle M_1K^{-1}M_1\rangle-\frac12{\dm}_{e^2}$. In momentum space the vertex of $M_1$ between $\phi(p)$ and $\bar\phi(q)$ is $e^{-\mathrm ik_\alpha/2}\sin\frac{(p+q)_\alpha}2$ ($k=q-p$), and $\langle A_\alpha^2\rangle=\intk D_b$; so $\Pi(p)=2\intk D_b(k)\Gh_a(p-k)\sum_\alpha\sin^2\frac{(2p-k)_\alpha}2-\frac12\intk D_b\sum_\alpha\cos p_\alpha-\frac12{\dm}_{e^2}$, which is $-\frac12\hat\Pi_{a,b}(p)-\frac12{\dm}_{e^2}$ (substitute $k\mapsto-k$ in $\hat\Pi$). By \cite[(3.33)]{K1} the full order-$e^2$ one-particle-irreducible two-point function vanishes at zero momentum, so $\Pi(p)=-\frac12\bar\pi_{a,b}(p)$. With $K^{-1}=2\Gh_a$, $\langle|\phi(0)|^2\rangle_{e^2}=\intp4\Gh_a^2\Pi=-2\intp\Gh_a^2\bar\pi_{a,b}=X_1(a,b)$.

(ii) Substitute $p=p'/\eps$, $k=k'/\eps$ in the same expression on $\eps\Z^3$: $\hat p_\eps=\eps^{-1}\hat p'$, $\Gh_\eps=\eps^2\Gh_{\eps m}$, $D_\eps=\eps^2D_{\eps\mu}$, the vertex scales as $\eps^{-1}$, so $\hat\Pi_\eps(p)=\eps^{-1}\hat\Pi_{\eps m,\eps\mu}(p')$ and $X_\eps=\eps^{-3}\eps^4\eps^{-1}X_1$.
\end{proof}

\begin{lemma}[a lattice identity]\label{lem:a05:trig}
For real $u,w$, with $\hat x^2\defeq4\sin^2(x/2)$: $\;4\sin^2\frac{u+w}2=2\hat u^2+2\hat w^2-\widehat{(u-w)}^2-\hat u^2\hat w^2$.
\end{lemma}
\begin{proof}
With $\hat x^2=2-2\cos x$ the right side is $2+2\cos(u-w)-4\cos u\cos w=2-2\cos u\cos w+2\sin u\sin w=2-2\cos(u+w)$.
\end{proof}

\begin{lemma}[exact decomposition of the order-$e^2$ self-energy]\label{lem:a05:dec}
For $a,b>0$ and $p\in\R^3$ put
\[
B_{a,b}(p)\defeq\intk D_b(k)\Gh_a(p-k),\qquad J^\alpha_{a,b}(p)\defeq\intk D_b(k)\Gh_a(p-k)\,\widehat{(p-k)}_\alpha^{\,2},\qquad d_b\defeq\intk D_b(k).
\]
Then
\[
-\tfrac12\bar\pi_{a,b}(p)=\hat p^2B_{a,b}(p)+\tfrac12(b^2-2a^2)\big(B_{a,b}(p)-B_{a,b}(0)\big)+\rho_{a,b}(p),\qquad
\rho_{a,b}(p)\defeq-\tfrac12\sum_\alpha\hat p_\alpha^2J^\alpha_{a,b}(p)+\tfrac14d_b\,\hat p^2,
\]
and therefore
\[
X_1(a,b)=4\intp\Gh_a(p)^2\Big[\hat p^2B_{a,b}(p)+\tfrac12(b^2-2a^2)\big(B_{a,b}(p)-B_{a,b}(0)\big)+\rho_{a,b}(p)\Big].
\]
\end{lemma}
\begin{proof}
Substitute $k\mapsto-k$ in $\hat\Pi$ ($D_b$ is even and the integrand is periodic): $\hat\Pi(p)=d_b\sum_\alpha\cos p_\alpha-\intk D_b(k)\Gh_a(p-k)\sum_\alpha4\sin^2\frac{(2p-k)_\alpha}2$. \Cref{lem:a05:trig} with $u=p_\alpha$, $w=(p-k)_\alpha$ ($u+w=(2p-k)_\alpha$, $u-w=k_\alpha$), summed over $\alpha$:
\[
\textstyle\sum_\alpha4\sin^2\frac{(2p-k)_\alpha}2=2\hat p^2+2\widehat{(p-k)}^2-\hat k^2-\sum_\alpha\hat p_\alpha^2\widehat{(p-k)}_\alpha^2
=2\hat p^2-2a^2+b^2+2\Gh_a(p-k)^{-1}-D_b(k)^{-1}-\sum_\alpha\hat p_\alpha^2\widehat{(p-k)}_\alpha^2 .
\]
Multiply by $D_b(k)\Gh_a(p-k)$ and integrate; $\intk\Gh_a(p-k)=\intk\Gh_a=:g_a$ by periodicity:
$\intk D_b\Gh_a(p-\cdot)\sum_\alpha4\sin^2\frac{(2p-k)_\alpha}2=(2\hat p^2-2a^2+b^2)B(p)+2d_b-g_a-\sum_\alpha\hat p_\alpha^2J^\alpha(p)$. With $\sum_\alpha\cos p_\alpha=3-\frac12\hat p^2$:
$\hat\Pi(p)=d_b-\frac12d_b\hat p^2-(2\hat p^2-2a^2+b^2)B(p)+g_a+\sum_\alpha\hat p_\alpha^2J^\alpha(p)$. Subtract the value at $p=0$, where $\hat p=0$.
\end{proof}

The same algebra with $(2p-k)^2=2p^2+2(p-k)^2-k^2$ gives the continuum function $\pi^{\rm c}(q)=(2q^2-2m^2+\mu^2)\Bc(q)-(\mu^2-2m^2)\Bc(0)$; \cref{lem:a05:dec} is its exact lattice form, the only lattice-specific terms being $\rho$ (of degree $2$ in $\hat p$) and the difference between the lattice and the continuum bubbles.

\begin{lemma}[elementary]\label{lem:a05:el}
\textup{(a)} For $|x|\le\pi$: $\frac4{\pi^2}x^2\le4\sin^2\frac x2\le x^2$ and $0\le x^2-4\sin^2\frac x2\le x^4/12$. For all real $x$: $4\sin^2\frac x2\ge\frac4{\pi^2}\dist(x,2\pi\Z)^2$. \textup{(b)} For $p\in\BZ$, $\dist(p,2\pi\Z^3)=|p|$. \textup{(c)} Hence for $q\in\BZ$: $\frac4{\pi^2}|q|^2\le\hat q^2\le|q|^2$, $0\le|q|^2-\hat q^2\le|q|^4/12$; and $\hat q^2\ge\frac4{\pi^2}\dist(q,2\pi\Z^3)^2$ for all $q\in\R^3$.
\end{lemma}
\begin{proof}
(a) $\sin y\ge\frac2\pi y$ on $[0,\frac\pi2]$ (concavity); $1-\frac{x^2}2\le\cos x\le1-\frac{x^2}2+\frac{x^4}{24}$ gives $x^2-\frac{x^4}{12}\le2-2\cos x\le x^2$; the last claim by periodicity. (b) $|p_\alpha-2\pi n_\alpha|\ge|p_\alpha|$ if $|p_\alpha|\le\pi$. (c) Sum over $\alpha$; $\sum_\alpha q_\alpha^4\le|q|^4$.
\end{proof}

\begin{lemma}[bubble comparison]\label{lem:a05:bub}
Let $\Bc_{a,b}(p)\defeq\int_{\R^3}\frac{d^3k}{(2\pi)^3}\frac1{(|k|^2+b^2)(|p-k|^2+a^2)}$. Then:
\begin{enumerate}[label=\textup{(\roman*)},nosep,leftmargin=*]
\item $\Bc_{a,b}(p)=\dfrac{\arctan(|p|/(a+b))}{4\pi|p|}$, $\Bc_{a,b}(0)=\dfrac1{4\pi(a+b)}$; $|p|\mapsto\Bc$ is non-increasing, $0<\Bc_{a,b}(p)\le\min\{\frac1{4\pi(a+b)},\frac1{8|p|}\}$, and $\Bc_{ta,tb}(tq)=t^{-1}\Bc_{a,b}(q)$.
\item For all $a,b>0$ and $p\in\BZ$: $|B_{a,b}(p)-\Bc_{a,b}(p)|\le C_0\defeq2\cdot10^3$.
\item For $p\in\BZ\setminus\{0\}$ and fixed $a,b>0$, as $t\downarrow0$: $B_{ta,tb}(p)\uparrow B_{0,0}(p)\defeq\intk\hat k^{-2}\widehat{(p-k)}^{-2}<\infty$ and $\Bc_{ta,tb}(p)\to\frac1{8|p|}$.
\end{enumerate}
\end{lemma}
\begin{proof}
(i) $\int_{\R^3}\frac{e^{\mathrm ik\cdot x}d^3k}{(2\pi)^3(|k|^2+b^2)}=\frac{e^{-b|x|}}{4\pi|x|}$, so $\Bc(p)=\int d^3x\,e^{-\mathrm ip\cdot x}\frac{e^{-(a+b)|x|}}{16\pi^2|x|^2}=\frac1{4\pi}\int_0^\infty e^{-(a+b)r}\frac{\sin(|p|r)}{|p|r}dr=\frac{\arctan(|p|/(a+b))}{4\pi|p|}$. $y\mapsto\arctan(y)/y$ is decreasing and $\le1$, and $\arctan\le\frac\pi2$. Homogeneity is read off.

(ii) \emph{Case $|p|\ge\frac12$.} $\Bc\le\frac14$ and $0<B_{a,b}(p)\le B_{0,0}(p)$. Split $\BZ$ into $V_1=\{|k|<\frac14\}$, $V_2=\{\dist(k-p,2\pi\Z^3)<\frac14\}$ (disjoint from $V_1$, because $\dist(p,2\pi\Z^3)=|p|\ge\frac12$ by \cref{lem:a05:el}(b)) and the rest. On $V_1$, $\dist(p-k,2\pi\Z^3)\ge\frac14$, so $\widehat{(p-k)}^{-2}\le4\pi^2$ (\cref{lem:a05:el}(c)), and $\int_{V_1}\hat k^{-2}\le\frac{\pi^2}4\int_{|k|<1/4}\frac{d^3k}{(2\pi)^3|k|^2}=\frac1{32}$. On $V_2$, $\dist(k,2\pi\Z^3)\ge\frac14$, so $\hat k^{-2}\le4\pi^2$, and by periodicity $\int_{V_2}\widehat{(p-k)}^{-2}\le\frac1{32}$. On the rest both factors are $\le4\pi^2$. So $B_{0,0}(p)\le\frac{\pi^2}4+16\pi^4<1.6\cdot10^3$.

\emph{Case $|p|<\frac12$.} Let $U=\{|k|<2\}\subset\BZ$. For $k\in U$, $k$ and $p-k$ lie in the ball of radius $\frac52<\pi$, which is inside $\BZ$. By \cref{lem:a05:el}(c), $0\le D_b(k)-(|k|^2+b^2)^{-1}=\frac{|k|^2-\hat k^2}{(\hat k^2+b^2)(|k|^2+b^2)}\le\frac{|k|^4/12}{(4/\pi^2)|k|^4}=\frac{\pi^2}{48}$, and the same for $\Gh_a(p-k)$. So on $U$, $|D_b(k)\Gh_a(p-k)-(|k|^2+b^2)^{-1}(|p-k|^2+a^2)^{-1}|\le\frac{\pi^2}{48}\big(\frac{\pi^2}4|p-k|^{-2}+|k|^{-2}\big)$, whose integral over $U$ is $\le\frac{\pi^2}{48}\big(\frac{\pi^2}4\cdot\frac5{4\pi^2}+\frac1{\pi^2}\big)<0.09$. On $\BZ\setminus U$: $\hat k^2\ge\frac{16}{\pi^2}$, and $\dist(p-k,2\pi\Z^3)\ge\frac32$ (for $n=0$, $|p-k|\ge\frac32$; for $n\ne0$ some $|p_\alpha-k_\alpha-2\pi n_\alpha|\ge2\pi-\frac12-\pi$), so the lattice integrand is $\le\frac{\pi^2}{16}\cdot\frac{\pi^2}9$ and its integral is $<0.7$. On $\R^3\setminus U$: $|p-k|\ge\frac34|k|$, so the continuum integrand is $\le\frac{16}9|k|^{-4}$, with integral $\frac4{9\pi^2}<0.05$. Total $<1$.

(iii) The integrand increases as $t\downarrow0$ to $\hat k^{-2}\widehat{(p-k)}^{-2}$, which is integrable for $p\notin2\pi\Z^3$ by the splitting of the first case with radius $|p|/4$; monotone convergence. The continuum limit is read off from (i).
\end{proof}

\begin{lemma}[$d$ and $J$]\label{lem:a05:J}
$d_b\le d_0\defeq\intk\hat k^{-2}\le\frac{\pi^2}4\int_\BZ\frac{d^3k}{(2\pi)^3|k|^2}<\infty$, $0\le J^\alpha_{a,b}(p)\le d_b$, hence $|\rho_{a,b}(p)|\le\frac34d_0\,\hat p^2$. As $t\downarrow0$, $d_{tb}\uparrow d_0$ and $J^\alpha_{ta,tb}(p)\to J^\alpha_{0,0}(p)\defeq\intk\hat k^{-2}\widehat{(p-k)}_\alpha^2\,\widehat{(p-k)}^{-2}$.
\end{lemma}
\begin{proof}
$\Gh_a(q)\hat q_\alpha^2\le1$; dominated convergence with the integrable majorant $\hat k^{-2}$ (the integrand converges off the null set $k\in p+2\pi\Z^3$).
\end{proof}

\begin{theorem}[asymptotics of the order-$e^2$ Wick constant]\label{thm:a05:Z}
Fix $m,\mu>0$. As $t\downarrow0$,
\[
X_1(tm,t\mu)=\frac1{4\pi^2}\ln\frac1t+\kappa(m,\mu)+o(1),\qquad\kappa(m,\mu)=\kappa_{\rm lat}+\kappa_{\rm c}(m,\mu),
\]
where, with $r_0(p)\defeq B_{0,0}(p)-\frac1{8|p|}$ and $\rho_0(p)\defeq-\frac12\sum_\alpha\hat p_\alpha^2J^\alpha_{0,0}(p)+\frac14d_0\hat p^2$,
\begin{align*}
\kappa_{\rm lat}&\defeq4\intp\frac{r_0(p)}{\hat p^2}+4\intp\frac{\rho_0(p)}{\hat p^4}+4\intp\Big(\frac1{\hat p^2}-\frac1{|p|^2}\Big)\frac1{8|p|}+\frac{\ln\pi}{4\pi^2}+\frac1{16\pi^3}\int_{\BZ\setminus\{|p|<\pi\}}\frac{d^3p}{|p|^3},\\
\kappa_{\rm c}(m,\mu)&\defeq-\frac{\ln m+\frac12}{4\pi^2}+\frac1{2\pi^3}\int_{\R^3}\frac{u_{m,\mu}(q)}{(q^2+m^2)^2}\,d^3q,\qquad
u_{m,\mu}(q)\defeq q^2\Bc_{m,\mu}(q)+\tfrac12(\mu^2-2m^2)\big(\Bc_{m,\mu}(q)-\Bc_{m,\mu}(0)\big)-\tfrac{|q|}8 .
\end{align*}
All integrals converge absolutely; $\kappa_{\rm lat}$ does not depend on $(m,\mu)$, and $\kappa_{\rm c}(m,\mu)=-\frac{\ln m}{4\pi^2}+\kappa_{\rm c}(1,\mu/m)$.
\end{theorem}
\begin{proof}
Put $a=tm$, $b=t\mu$, $t\in(0,1]$, $m_*\defeq\max(m,\mu)$, $r\defeq B-\Bc$ ($|r|\le C_0$ on $\BZ$, \cref{lem:a05:bub}(ii)), and $G^{\rm c}_a(p)\defeq(|p|^2+a^2)^{-1}$. Two integrals are used repeatedly: $\intp\hat p^{-2}<\infty$, and
\[
\intp\Gh_a^2\le\int_{\R^3}\frac{d^3p}{(2\pi)^3((4/\pi^2)p^2+a^2)^2}=\frac{\pi^2}{64a}.
\]
By \cref{lem:a05:dec}, $X_1=M(t)+E_1+E_2+E_3$ with
\begin{gather*}
M(t)\defeq4\intp\Gh_a^2\big[\hat p^2\Bc+\tfrac12(b^2-2a^2)(\Bc(p)-\Bc(0))\big],\qquad E_1\defeq4\intp\Gh_a^2\hat p^2r,\\
E_2\defeq2(b^2-2a^2)\intp\Gh_a^2(r(p)-r(0)),\qquad E_3\defeq4\intp\Gh_a^2\rho .
\end{gather*}
\emph{Step 1.} $|4\Gh_a^2\hat p^2r|\le4C_0\hat p^{-2}$ and $|4\Gh_a^2\rho|\le3d_0\hat p^{-2}$; the integrands converge pointwise on $\BZ\setminus\{0\}$ (\cref{lem:a05:bub}(iii), \cref{lem:a05:J}; $\Gh_{ta}(p)\to\hat p^{-2}$). By dominated convergence $E_1\to4\intp r_0\hat p^{-2}$ and $E_3\to4\intp\rho_0\hat p^{-4}$.

\emph{Step 2.} $|E_2|\le6m_*^2t^2\cdot2C_0\cdot\frac{\pi^2}{64tm}=O(t)$.

\emph{Step 3.} Write $M(t)=M^{\rm c}(t)+E_4+E_5$, where $M^{\rm c}$ is $M$ with $\Gh_a^2\hat p^2$ replaced by $(G^{\rm c}_a)^2|p|^2$ and $\Gh_a^2$ by $(G^{\rm c}_a)^2$ (still integrated over the cube $\BZ$), $E_4\defeq4\intp\big(\Gh_a^2\hat p^2-(G^{\rm c}_a)^2|p|^2\big)\Bc$, and $E_5\defeq2(b^2-2a^2)\intp\big(\Gh_a^2-(G^{\rm c}_a)^2\big)(\Bc(p)-\Bc(0))$. For $f(x)=x(x+a^2)^{-2}$, $|f'(x)|\le(x+a^2)^{-2}$; both $\hat p^2$ and $|p|^2$ lie in $[\frac4{\pi^2}|p|^2,|p|^2]$ and differ by $\le|p|^4/12$, so $|\Gh_a^2\hat p^2-(G^{\rm c}_a)^2|p|^2|\le\frac{\pi^4}{192}$. With $\Bc\le\frac1{8|p|}$ the integrand of $E_4$ has an integrable majorant, and it converges pointwise; so $E_4\to4\intp(\hat p^{-2}-|p|^{-2})\frac1{8|p|}$. For $g(x)=(x+a^2)^{-2}$, $|g'(x)|\le2(x+a^2)^{-3}$, hence $|\Gh_a^2-(G^{\rm c}_a)^2|\le\frac{\pi^4}{96}\big(\frac4{\pi^2}|p|^2+a^2\big)^{-1}$; with $|\Bc(p)-\Bc(0)|\le\frac1{4\pi t(m+\mu)}$, $|E_5|\le6m_*^2t^2\cdot\frac{\pi^4}{96}\cdot\frac1{4\pi t(m+\mu)}\cdot\frac{\pi^2}4\intp|p|^{-2}=O(t)$.

\emph{Step 4: $M^{\rm c}(t)$.} Fix $\varrho\in(0,1]$ and split $\BZ$ into the ball $\{|p|<\varrho\}$ and the rest. On the rest the first part of the integrand, $4(G^{\rm c}_a)^2|p|^2\Bc$, is $\le\frac1{2|p|^3}$ and converges to $\frac1{2|p|^3}$; the second part is $\le|p|^{-4}\cdot3m_*^2t^2\cdot\frac1{4\pi t(m+\mu)}\to0$ uniformly on $|p|\ge\varrho$; so the outer part tends to $\int_{\BZ\setminus\{|p|<\varrho\}}\frac{d^3p}{(2\pi)^3\,2|p|^3}=\frac{\ln(\pi/\varrho)}{4\pi^2}+\frac1{16\pi^3}\int_{\BZ\setminus\{|p|<\pi\}}\frac{d^3p}{|p|^3}$. On the ball substitute $p=tq$; by the homogeneity in \cref{lem:a05:bub}(i) and $G^{\rm c}_{ta}(tq)=t^{-2}G^{\rm c}_a(q)$, the inner part equals exactly
\[
\frac4{(2\pi)^3}\int_{|q|<\varrho/t}\frac{q^2\Bc_{m,\mu}(q)+\frac12(\mu^2-2m^2)(\Bc_{m,\mu}(q)-\Bc_{m,\mu}(0))}{(q^2+m^2)^2}\,d^3q
=\frac1{2\pi^3}\int_{|q|<\varrho/t}\frac{|q|/8+u_{m,\mu}(q)}{(q^2+m^2)^2}\,d^3q .
\]
$u_{m,\mu}$ is continuous and bounded on $\R^3$ ($u(0)=0$; $u(q)\to-\frac{m+\mu}{4\pi}-\frac{\mu^2-2m^2}{8\pi(m+\mu)}$ as $|q|\to\infty$, from $\arctan y=\frac\pi2-\frac1y+O(y^{-3})$), so its part converges to $\frac1{2\pi^3}\int_{\R^3}u(q^2+m^2)^{-2}$. The $|q|/8$ part is $\frac1{4\pi^2}\int_0^Q\frac{q^3dq}{(q^2+m^2)^2}=\frac1{8\pi^2}\big[\ln\frac{Q^2+m^2}{m^2}+\frac{m^2}{Q^2+m^2}-1\big]=\frac1{4\pi^2}\big(\ln\frac Qm-\frac12\big)+o(1)$, $Q=\varrho/t$. Adding inner and outer parts, $\ln\varrho$ cancels and $M^{\rm c}(t)=\frac1{4\pi^2}\ln\frac1t+\kappa_{\rm c}(m,\mu)+\frac{\ln\pi}{4\pi^2}+\frac1{16\pi^3}\int_{\BZ\setminus\{|p|<\pi\}}|p|^{-3}d^3p+o(1)$.

Collecting Steps 1--4 gives the claim. Absolute convergence of the integrals in $\kappa_{\rm lat}$ is the majorant of Steps 1 and 3. The last statement is the substitution $q=mq'$ and the degree-one homogeneity of $u$ in $(q,m,\mu)$.
\end{proof}

\begin{proof}[Proof of \cref{lem:k02:cnum}]
(a) By \cref{lem:a05:deriv}(ii) and \cref{thm:a05:Z} with $t=\eps_k=5^{-k}$, $z_k=e^2X_1(\eps_km,\eps_k\lmu)=\frac{e^2\ln5}{4\pi^2}\,k+e^2\kappa(m,\lmu)+o(1)$, so $|z_k|\le C(1+k)$ with $C$ depending on the couplings only. (Also $z_k\to+\infty$, and $z_k-z_{k-1}\to\frac{e^2}{4\pi^2}\ln5$.)

(b) Fix $M>0$ and write $\sW_\eps(M)=\int_{[-\pi/\eps,\pi/\eps)^3}\frac{d^3q}{(2\pi)^3}\,(\hat q^2+M^2)^{-1}$, $\hat q^2=\sum_\alpha\frac4{\eps^2}\sin^2\frac{\eps q_\alpha}2$. Then $\sW_\eps(m)-\sW_\eps(M)=(M^2-m^2)\int\frac{d^3q}{(2\pi)^3}\frac1{(\hat q^2+m^2)(\hat q^2+M^2)}$; by \cref{lem:a05:el}(c) (scaled to $\eps\Z^3$) the integrand is bounded by $\frac{\pi^4}{16}\big((q^2+m^2)(q^2+M^2)\big)^{-1}$, integrable on $\R^3$, and converges pointwise, so the Wick difference converges to $(M^2-m^2)\frac1{4\pi(m+M)}=\frac{M-m}{4\pi}$ and is bounded in $k$. For the scheme difference, \cref{thm:a05:Z} at the masses $(m,\lmu)$ and $(M,\lmu)$ with the same $t=\eps_k$ gives $z_k(m)-z_k(M)=e^2\big[X_1(\eps_km,\eps_k\lmu)-X_1(\eps_kM,\eps_k\lmu)\big]\to e^2\big[\kappa_{\rm c}(m,\lmu)-\kappa_{\rm c}(M,\lmu)\big]$: the logarithms and $\kappa_{\rm lat}$ cancel. Hence $\Delta_k(M)$ converges, in particular $\sup_k|\Delta_k(M)|<\infty$.
\end{proof}

\subsection{King's bare coefficient and its growth with the mass}\label{app:a05:B}

\paragraph{The object.} King's bare coefficient at scalar mass $M$ and spacing $\eps$ is $\alpha_\eps(M)=M^2+{\dm}_\eps(M)$, where ${\dm}_\eps(M)$ is defined by \cite[(3.32)--(3.33)]{K1} with the propagators of the infinite lattice $\eps\Z^3$ \cite[p.~330]{K2}: it is the finite sum
\[
{\dm}_\eps(M)=\lambda\,d_{01}+e^2d_{20}+e^4d_{40}+\lambda e^2d_{21}+\lambda^2d_{02},
\]
over the orders $2\le\alpha+2\beta\le4$ in $e^\alpha\lambda^\beta$, each coefficient $d_{\alpha\beta}=d_{\alpha\beta}(M^2;\eps)$ being the zero-momentum value of the one-particle-irreducible two-point function of its order with all lower-order counterterms inserted; the odd powers of $e$ vanish, since the action is invariant under $(A,\phi)\mapsto(-A,R\phi)$ with $R$ a reflection of the two components. There is no photon and no wave-function counterterm. With $\alpha_k(M)\defeq\alpha_{\eps_k}(M)$, the quantities of \cref{ssec:k02:setting} are $u_k(M)=\frac12(\alpha_k(M)-\alpha_k(m))$ and $s_-(M)=\liminf_ku_k(M)$.

\paragraph{Lattice formulas.} On $\eps\Z^3$ put $\BZ_\eps\defeq[-\pi/\eps,\pi/\eps)^3$, $\langle f\rangle\defeq\int_{\BZ_\eps}f\,\frac{d^3q}{(2\pi)^3}$, $\langle\langle f\rangle\rangle$ for the double integral, $\hat q^2\defeq\sum_\alpha\frac4{\eps^2}\sin^2\frac{\eps q_\alpha}2$, $s\defeq M^2$, $C(q)\defeq(\hat q^2+s)^{-1}$ and $D(k)\defeq(\hat k^2+\lmu^2)^{-1}$. Evaluating King's graphs \cite[(3.34)]{K1} with the Feynman rules of his action (the current vertex of \cref{lem:a05:deriv}, the seagull, the quartic vertex $\lambda|\phi|^4$) gives, with $y_{\alpha\beta}\defeq\partial_sd_{\alpha\beta}$:
\begin{enumerate}[label=\textup{(\roman*)},nosep,leftmargin=*]
\item $d_{01}=-16\langle C\rangle$ (the quartic tadpole), $y_{01}=16\langle C^2\rangle>0$;
\item $d_{20}=\langle\hat q^2DC\rangle-3\langle D\rangle$ (photon exchange and seagull tadpole), $y_{20}=-\langle\hat q^2DC^2\rangle<0$;
\item $d_{02}=128\langle\langle C(q)C(r)C(q+r)\rangle\rangle$ (the sunset; the tadpole-on-tadpole cancels identically against the inserted $d_{01}$), $y_{02}=-384\langle\langle C(q)^2C(r)C(q+r)\rangle\rangle<0$;
\item $d_{21}=-8X_\eps(M;\lmu)=-16\langle C^2\sigma_\eps\rangle$, where $\sigma_\eps\defeq\Sigma^{(20)}_{\rm ren}/e^2$ is the renormalised order-$e^2$ self-energy per real component ($-\bar\pi$ of \cref{lem:a05:deriv} scaled to $\eps\Z^3$, twice the function $-\frac12\bar\pi$ of \cref{lem:a05:dec}; cf.\ \eqref{eq:a05:sigma}): the quartic tadpole with the order-$e^2$ self-energy inserted on its line is the only non-vanishing graph of this order, the tadpole-on-tadpole being cancelled by the inserted $d_{01}$ and the two remaining classes vanishing identically at zero external momentum by the oddness of the current vertex;
\item $d_{40}$ is the sum over the twelve two-loop topologies of order $e^4$ (listed in the proof of \cref{lem:a05:D45}(v) and, by continuum class, in \cref{lem:a05:40} below), with the order-$e^2$ counterterm inserted on the exchange line.
\end{enumerate}
Each $d_{\alpha\beta}(s;\eps)$ is a smooth function of $s>0$ at fixed $\eps$ (an integral of a smooth integrand over a compact zone), and the \emph{exact identity}
\begin{equation}\label{eq:a05:exact}
\alpha_\eps(M)-\alpha_\eps(m)=(M^2-m^2)+\sum_{\alpha\beta}e^\alpha\lambda^\beta\int_{m^2}^{M^2}y_{\alpha\beta}(s;\eps)\,ds
\end{equation}
holds for all $M,m>0$. The formula (i), (ii) at one loop is the one behind \eqref{eq:s14:alphaprime}.

\begin{lemma}[one-loop and sunset bounds on $\eps\Z^3$]\label{lem:a05:D123}
For every $\eps>0$ and $s>0$:
\textup{(i)} $0<y_{01}(s)\le\dfrac{\pi^2}{4\sqrt s}$;
\textup{(ii)} $-\dfrac{\pi^2}{32(\sqrt s+\lmu)}\le y_{20}(s)<0$;
\textup{(iii)} $-\dfrac{81\pi^4}{32\,s}\le y_{02}(s)<0$.
\end{lemma}
\begin{proof}
By \cref{lem:a05:el}(c) scaled to $\eps\Z^3$, $\hat q^2\ge\frac4{\pi^2}|q|^2$ on $\BZ_\eps$, so $C(q)\le\frac{\pi^2}4\tilde G(q)$ with $\tilde G(q)\defeq(|q|^2+\tilde a^2)^{-1}$, $\tilde a\defeq\frac\pi2\sqrt s$, and $D(k)\le\frac{\pi^2}4(|k|^2+\tilde b^2)^{-1}$, $\tilde b\defeq\frac\pi2\lmu$; moreover $C(q+r)\le\frac{\pi^2}4\tilde G(q+r-P)$ for the $P\in\frac{2\pi}\eps\Z^3$ with $q+r-P\in\BZ_\eps$, and for $q,r\in\BZ_\eps$ there are $27$ candidates $P$.
(i) $16\langle C^2\rangle\le16\cdot\frac{\pi^4}{16}\int_{\R^3}\frac{d^3q}{(2\pi)^3}\tilde G^2=\pi^4\cdot\frac1{8\pi\tilde a}=\frac{\pi^2}{4\sqrt s}$.
(ii) $\hat q^2C\le1$ gives $|y_{20}|\le\langle DC\rangle\le\frac{\pi^4}{16}\int_{\R^3}\frac{d^3k}{(2\pi)^3}\frac1{(|k|^2+\tilde b^2)(|k|^2+\tilde a^2)}=\frac{\pi^4}{16}\cdot\frac1{4\pi(\tilde a+\tilde b)}=\frac{\pi^2}{32(\sqrt s+\lmu)}$ (\cref{lem:a05:bub}(i)); strictness, since every summand is positive.
(iii) $\langle\langle C(q)^2C(r)C(q+r)\rangle\rangle\le(\frac{\pi^2}4)^4\int\tilde G(q)^2\Big[\sum_P\int_{\R^3}\tilde G(r)\tilde G(q+r-P)\,\frac{d^3r}{(2\pi)^3}\Big]\frac{d^3q}{(2\pi)^3}\le(\frac{\pi^2}4)^4\cdot\frac1{8\pi\tilde a}\cdot\frac{27}{8\pi\tilde a}=\frac{27\pi^4}{4096\,s}$, by \cref{lem:a05:bub}(i) ($\Bc_{\tilde a,\tilde a}\le\frac1{8\pi\tilde a}$). Multiply by $384$.
\end{proof}

\begin{lemma}[two-loop lattice bounds on $\eps\Z^3$]\label{lem:a05:D45}
\textup{(iv)} For every $\eps>0$ and $s>0$, $|y_{21}(s;\eps)|\le C_{21}(s)/s+C'_{21}\,\eps/\sqrt s$, with
\[
C_{21}(s)\defeq\frac{\pi^2}4\Big[\frac{135}8\pi^2+\frac{135\pi^4}{128}\Big(2+\frac{\lmu^2}s\Big)+\frac{\pi^2}{16}\Big]\quad(\text{non-increasing in }s),\qquad C'_{21}\defeq\frac{3\sqrt3\,\pi^3}{32}.
\]
\textup{(v)} There are $a_0\in(0,1]$ and $C_{40}<\infty$, independent of $\lmu$, $\eps$ and $s$, such that $|y_{40}(s;\eps)|\le C_{40}/s$ whenever $\eps\le a_0/\sqrt s$.
\end{lemma}
% a05_twoloop_iv: proof of lem:a05:D45 (iv), the |y_21| bound on eps Z^3 (input from a05_cnum_NEG.tex)
\paragraph{Proof of \cref{lem:a05:D45}(iv).} The lattice is $\eps\Z^3$ with the conventions of the paragraph ``Lattice formulas'': $s=M^2$, $C(q)=(\hat q^2+s)^{-1}$, $D(k)=(\hat k^2+\lmu^2)^{-1}$, $\langle\cdot\rangle$ the integral over $\BZ_\eps$, $\bar\BZ_\eps$ the closed zone, $\hat x_\mu\defeq\frac2\eps\sin\frac{\eps x_\mu}2$. Every integrand below is smooth and $\frac{2\pi}\eps\Z^3$-periodic in each integration variable, so $k\mapsto p-k$ is a permitted substitution for every real $p$. Put
\begin{gather*}
B_s(p)\defeq\langle D(k)C(p-k)\rangle_k,\qquad \Xi_s(p)\defeq\langle D(k)C(p-k)^2\rangle_k=-\partial_sB_s(p),\qquad J^\alpha_s(p)\defeq\langle D(k)C(p-k)\widehat{(p-k)}_\alpha^{\,2}\rangle_k,\\
s_0\defeq\langle D\rangle,\qquad b_s\defeq\langle C^2\rangle,\qquad F_s(p)\defeq\langle\widehat{(p+r)}^2D(p-r)C(r)\rangle_r,\qquad G_s(p)\defeq\partial_sF_s(p)=-\langle\widehat{(p+r)}^2D(p-r)C(r)^2\rangle_r .
\end{gather*}
$F_s$ is the photon-exchange graph and $s_0\sum_\mu\cos\eps p_\mu$ the seagull tadpole at external momentum $p$: the function $\hat\Pi$ of \cref{lem:a05:deriv}, at spacing $\eps$ and after the substitution $r=p+k$, is $s_0\sum_\mu\cos\eps p_\mu-F_s(p)$. Hence the self-energy per real component is
\begin{equation}\label{eq:a05:sigeps}
\sigma_\eps(p)=-\bar\pi_\eps(p)=F_s(p)-F_s(0)+s_0\sum_\mu(1-\cos\eps p_\mu)=F_s(p)-F_s(0)+\tfrac12\eps^2s_0\,\hat p^2 ,
\end{equation}
and since $s_0$ does not depend on $s$, $\partial_s\sigma_\eps(q)=G_s(q)-G_s(0)$. Differentiating (iv) of the lattice formulas,
\begin{equation}\label{eq:a05:y21}
y_{21}=\partial_sd_{21}=-16\big\langle-2C^3\sigma_\eps+C^2\,\partial_s\sigma_\eps\big\rangle .
\end{equation}

\begin{lemma}[elementary lattice inequalities, continued]\label{lem:a05:el2}
Let $u,v\in\R^3$, $\tilde a\defeq\frac\pi2\sqrt s$, $\tilde b\defeq\frac\pi2\lmu$, $\tilde G_a(w)\defeq(|w|^2+a^2)^{-1}$ on $\R^3$.
\begin{enumerate}[label=\textup{(\alph*)},nosep,leftmargin=*]
\item $|\hat u_\mu-\hat v_\mu|\le|u_\mu-v_\mu|$, $|\widehat{(u+v)}_\mu|\le|\hat u_\mu|+|\hat v_\mu|$, $|\widehat{(2v)}_\mu|\le2|\hat v_\mu|$; hence $\widehat{(u+v)}^2\le2\,\widehat{(u-v)}^2+8\hat v^2$.
\item With $E(v)\defeq\hat v^2$: $\partial_\mu E=\frac2\eps\sin\eps v_\mu$, $|\partial_\mu E|\le2|\hat v_\mu|$, $\partial_\mu\partial_\nu E=2\cos(\eps v_\mu)\delta_{\mu\nu}$; and $|w\cdot\operatorname{Hess}C(v)\,w|\le10\,C(v)^2|w|^2$ for all $w\in\R^3$.
\item For $p\in\bar\BZ_\eps$ and $k\in\BZ_\eps$: $C(p-k)\le\frac{\pi^2}4\sum_{n\in\{-1,0,1\}^3}\tilde G_{\tilde a}\big(p-k-\tfrac{2\pi}\eps n\big)$ and $C(p-k)^2\le\big(\frac{\pi^2}4\big)^2\sum_n\tilde G_{\tilde a}\big(p-k-\tfrac{2\pi}\eps n\big)^2$; for $k\in\BZ_\eps$, $C(k)\le\frac{\pi^2}4\tilde G_{\tilde a}(k)$ and $D(k)\le\frac{\pi^2}4\tilde G_{\tilde b}(k)$.
\item $s_0\le d_0/\eps$ with $d_0=\intk\hat k^{-2}\le\frac{\sqrt3\,\pi}8$ ($d_0$ as in \cref{lem:a05:J}); $0\le J^\alpha_s(p)\le s_0$.
\item $\int_{\R^3}\tilde G_b(k)\tilde G_a(P-k)^2\frac{d^3k}{(2\pi)^3}=\dfrac1{8\pi a\,((a+b)^2+|P|^2)}\le\dfrac1{8\pi a(a+b)^2}$.
\end{enumerate}
\end{lemma}
\begin{proof}
(a) The derivative of $x\mapsto\frac2\eps\sin\frac{\eps x}2$ is $\cos\frac{\eps x}2$, of modulus $\le1$; the addition theorem gives the second claim and $|\sin2y|\le2|\sin y|$ the third. Then $|\widehat{(u+v)}_\mu|=|\widehat{((u-v)+2v)}_\mu|\le|\widehat{(u-v)}_\mu|+2|\hat v_\mu|$ and $(x+y)^2\le2x^2+2y^2$; sum over $\mu$.
(b) $|\frac2\eps\sin\eps v_\mu|=\frac4\eps|\sin\frac{\eps v_\mu}2\cos\frac{\eps v_\mu}2|\le2|\hat v_\mu|$. With $\partial_\mu C=-C^2\partial_\mu E$: $w\cdot\operatorname{Hess}C\,w=2C^3(w\cdot\nabla E)^2-2C^2\sum_\mu\cos(\eps v_\mu)w_\mu^2$, and $|w\cdot\nabla E|\le\sum_\mu2|w_\mu||\hat v_\mu|\le2|w||\hat v|$, so the modulus is $\le(8C^3\hat v^2+2C^2)|w|^2\le10C^2|w|^2$, since $C\hat v^2\le1$.
(c) The components of $p-k$ lie in $[-\frac{2\pi}\eps,\frac{2\pi}\eps]$, so a nearest point $\bar n$ of $\frac{2\pi}\eps\Z^3$ lies in $\frac{2\pi}\eps\{-1,0,1\}^3$. By \cref{lem:a05:el}(c) at spacing $\eps$, $\widehat{(p-k)}^2\ge\frac4{\pi^2}|p-k-\bar n|^2$, so $C(p-k)\le(\frac4{\pi^2}|p-k-\bar n|^2+s)^{-1}=\frac{\pi^2}4\tilde G_{\tilde a}(p-k-\bar n)$; a maximum of positive terms is at most their sum. For $k\in\BZ_\eps$, $\bar n=0$.
(d) $k=k'/\eps$ gives $\langle\hat k^{-2}\rangle=d_0/\eps$, and $d_0\le\frac{\pi^2}4\int_{|k|\le\sqrt3\pi}\frac{d^3k}{(2\pi)^3|k|^2}=\frac{\pi^2}4\cdot\frac{\sqrt3}{2\pi}$. $J^\alpha_s\le s_0$ because $C(v)\hat v_\alpha^2\le C(v)\hat v^2\le1$.
(e) By \cref{lem:a05:bub}(i) the integral with $\tilde G_a$ in place of $\tilde G_a^2$ is $\frac{\arctan(|P|/(a+b))}{4\pi|P|}$; apply $-\partial/\partial(a^2)=-(2a)^{-1}\partial_a$ and $\partial_a\arctan\frac x{a+b}=-\frac x{(a+b)^2+x^2}$.
\end{proof}

\begin{lemma}[the decomposition on $\eps\Z^3$]\label{lem:a05:Kdec}
For every $p\in\R^3$,
\[
F_s(p)-F_s(0)=2\hat p^2B_s(p)+(\lmu^2-2s)\big(B_s(p)-B_s(0)\big)-\eps^2\sum_\alpha\hat p_\alpha^2J^\alpha_s(p),
\]
hence $\sigma_\eps(p)=2\hat p^2B_s(p)+(\lmu^2-2s)(B_s(p)-B_s(0))+2\rho_\eps(p)$ with $\rho_\eps(p)\defeq-\frac12\eps^2\sum_\alpha\hat p_\alpha^2J^\alpha_s(p)+\frac14\eps^2s_0\hat p^2$. (This is \cref{lem:a05:dec} at spacing $\eps$.)
\end{lemma}
\begin{proof}
\Cref{lem:a05:trig} with $u=\eps x$, $w=\eps y$, divided by $\eps^2$, reads $\widehat{(x+y)}^2=2\hat x^2+2\hat y^2-\eps^2\hat x^2\hat y^2-\widehat{(x-y)}^2$ componentwise. With $x=p_\alpha$, $y=(p-k)_\alpha$, summed over $\alpha$: $\widehat{(2p-k)}^2=2\hat p^2+2\widehat{(p-k)}^2-\eps^2\sum_\alpha\hat p_\alpha^2\widehat{(p-k)}_\alpha^2-\hat k^2$. Substitute $r=p-k$ in $F_s$ and insert $\widehat{(p-k)}^2=C(p-k)^{-1}-s$, $\hat k^2=D(k)^{-1}-\lmu^2$:
\[
F_s(p)=\big\langle D(k)C(p-k)\big[2\hat p^2-2s+\lmu^2\big]\big\rangle_k+2s_0-\langle C\rangle-\eps^2\sum_\alpha\hat p_\alpha^2J^\alpha_s(p),
\]
using $\langle C(p-k)\rangle_k=\langle C\rangle$. Subtract the value at $p=0$, where $\hat p=0$: the linearly divergent constants $2s_0-\langle C\rangle$ cancel exactly. The second form follows from \eqref{eq:a05:sigeps}.
\end{proof}

\begin{lemma}[one-loop blocks with explicit $s$]\label{lem:a05:P}
For every $\eps>0$ and $s>0$:
\begin{enumerate}[label=\textup{(P\arabic*)},nosep,leftmargin=*]
\item $b_s\le\dfrac{\pi^2}{64\sqrt s}$;
\item $B_s(0)\le\dfrac{\pi^2}{32(\sqrt s+\lmu)}$ and $\sup_{p\in\bar\BZ_\eps}B_s(p)\le\dfrac{27\pi^2}{32(\sqrt s+\lmu)}$;
\item $\sup_{p\in\bar\BZ_\eps}\Xi_s(p)\le\dfrac{27\pi^2}{64\,s^{3/2}}$;
\item $|G_s(p)|\le2b_s+8B_s(p)$ for $p\in\bar\BZ_\eps$; hence $|\partial_s\sigma_\eps(q)|\le4b_s+16\sup_{\bar\BZ_\eps}B_s$ for $q\in\BZ_\eps$;
\item $|\sigma_\eps(p)|\le2\hat p^2B_s(p)+5(\lmu^2+2s)\,|p|^2\sup_{\bar\BZ_\eps}\Xi_s+\frac32\eps^2s_0\hat p^2$ for $p\in\BZ_\eps$.
\end{enumerate}
\end{lemma}
\begin{proof}
(P1) By \cref{lem:a05:el2}(c), $b_s\le(\frac{\pi^2}4)^2\int_{\R^3}\tilde G_{\tilde a}^2\frac{d^3q}{(2\pi)^3}=\frac{\pi^4}{16}\cdot\frac1{8\pi\tilde a}=\frac{\pi^2}{64\sqrt s}$.
(P2) At $p=0$ both factors are at $k\in\BZ_\eps$: $B_s(0)\le(\frac{\pi^2}4)^2\int\tilde G_{\tilde b}\tilde G_{\tilde a}=\frac{\pi^4}{16}\cdot\frac1{4\pi(\tilde a+\tilde b)}$ by \cref{lem:a05:bub}(i). At general $p\in\bar\BZ_\eps$, \cref{lem:a05:el2}(c) supplies $27$ images, each integral being $\le\frac1{4\pi(\tilde a+\tilde b)}$.
(P3) \Cref{lem:a05:el2}(c), (e): $\Xi_s(p)\le(\frac{\pi^2}4)^3\sum_n\int\tilde G_{\tilde b}(k)\tilde G_{\tilde a}(p-k-\tfrac{2\pi}\eps n)^2\frac{d^3k}{(2\pi)^3}\le\frac{27(\pi^2/4)^3}{8\pi\tilde a(\tilde a+\tilde b)^2}\le\frac{27\pi^6/64}{8\pi(\pi/2)^3s^{3/2}}=\frac{27\pi^2}{64\,s^{3/2}}$.
(P4) \Cref{lem:a05:el2}(a) with $u=p$, $v=r$: $\widehat{(p+r)}^2D(p-r)C(r)^2\le\big[2\widehat{(p-r)}^2D(p-r)+8\hat r^2D(p-r)\big]C(r)^2\le2C(r)^2+8D(p-r)C(r)$, by $\hat x^2(\hat x^2+\lmu^2)^{-1}\le1$ and $\hat r^2C(r)\le1$. Integrate over $r$; $\langle D(p-r)C(r)\rangle_r=B_s(p)$. The second claim is $\partial_s\sigma_\eps=G_s(q)-G_s(0)$.
(P5) In \cref{lem:a05:Kdec}, $\eps^2\sum_\alpha\hat p_\alpha^2J^\alpha_s\le\eps^2\hat p^2s_0$ by \cref{lem:a05:el2}(d), and the seagull term is $\frac12\eps^2s_0\hat p^2$. For the middle term: $B_s$ is smooth on $\R^3$ and even ($k\mapsto-k$; $D$, $C$ even), so $\nabla B_s(0)=0$ and $B_s(p)-B_s(0)=\int_0^1(1-t)\,(p\cdot\nabla)^2B_s(tp)\,dt$ with $(p\cdot\nabla)^2B_s(tp)=\langle D(k)\,p\cdot\operatorname{Hess}C(tp-k)\,p\rangle_k$. By \cref{lem:a05:el2}(b), $|(p\cdot\nabla)^2B_s(tp)|\le10|p|^2\Xi_s(tp)$ with $tp\in\bar\BZ_\eps$, and $\int_0^1(1-t)dt=\frac12$. Finally $|\lmu^2-2s|\le\lmu^2+2s$.
\end{proof}

\begin{proof}[Proof of (iv)]
By \eqref{eq:a05:y21}, $|y_{21}|\le16\big[2\langle C^3|\sigma_\eps|\rangle+\langle C^2|\partial_s\sigma_\eps|\rangle\big]$. For $q\in\BZ_\eps$, $C^3\hat q^2\le C^2$ and $C^3|q|^2\le\frac{\pi^2}4C^3\hat q^2\le\frac{\pi^2}4C^2$ (\cref{lem:a05:el}(c)). So \cref{lem:a05:P}(P5) gives $2\langle C^3|\sigma_\eps|\rangle\le b_s\big[4\sup B_s+\frac{5\pi^2}2(\lmu^2+2s)\sup\Xi_s+3\eps^2s_0\big]$, and (P4) gives $\langle C^2|\partial_s\sigma_\eps|\rangle\le b_s\big[4b_s+16\sup B_s\big]$, the suprema over $\bar\BZ_\eps$. Hence
\[
|y_{21}|\le16\,b_s\Big[20\sup B_s+\frac{5\pi^2}2(\lmu^2+2s)\sup\Xi_s+4b_s+3\eps^2s_0\Big].
\]
Insert (P1)--(P3) and \cref{lem:a05:el2}(d), dropping $\lmu\ge0$ from the denominators: $16b_s\le\frac{\pi^2}{4\sqrt s}$; $20\sup B_s\le\frac{135}8\frac{\pi^2}{\sqrt s}$; $\frac{5\pi^2}2(\lmu^2+2s)\cdot\frac{27\pi^2}{64s^{3/2}}=\frac{135\pi^4}{128}\big(2+\frac{\lmu^2}s\big)\frac1{\sqrt s}$; $4b_s\le\frac{\pi^2}{16\sqrt s}$; $3\eps^2s_0\le\frac{3\sqrt3\,\pi}8\eps$. Collecting,
\[
|y_{21}(s;\eps)|\le\frac{\pi^2}4\Big[\frac{135}8\pi^2+\frac{135\pi^4}{128}\Big(2+\frac{\lmu^2}s\Big)+\frac{\pi^2}{16}\Big]\frac1s+\frac{\pi^2}4\cdot\frac{3\sqrt3\,\pi}8\cdot\frac\eps{\sqrt s}=\frac{C_{21}(s)}s+C'_{21}\frac\eps{\sqrt s},
\]
and $C_{21}$ is non-increasing in $s$ because $\lmu^2/s$ is. The only coefficient that grows with $s$, $|\lmu^2-2s|$, has been paired with $B_s(p)-B_s(0)=O(|p|^2s^{-3/2})$, the Taylor step being taken on the scalar line.
\end{proof}

% a05_twoloop_v: proof of lem:a05:D45 (v), part 1: scaling, forest form, the twelve graphs, lemmas (input from a05_cnum_NEG.tex)
\paragraph{Proof of \cref{lem:a05:D45}(v).} The proof rescales to unit scalar mass, writes King's coefficient as a forest-subtracted sum over the twelve graphs of order $e^4$, and bounds that sum uniformly in the rescaled photon mass. The power-counting input of \cite{Rei88,Rei88b} enters at one mass configuration only (scalar and photon mass $1$); the light-photon limit is handled by elementary bounds.

\begin{lemma}[exact scaling]\label{lem:a05:SC}
Write $d_{\alpha\beta}(s;\lmu;\eps)$ for the coefficient at renormalised scalar mass squared $s$, photon mass $\lmu$ and spacing $\eps$, and $n\defeq\alpha/2+\beta$. For every $t>0$, $d_{\alpha\beta}(s;\lmu;\eps)=t^{2-n}\,d_{\alpha\beta}(s/t^2;\lmu/t;t\eps)$. Consequently
\[
y_{40}(s;\lmu;\eps)=s^{-1}\,Y\big(\lmu/\sqrt s,\ \eps\sqrt s\big),\qquad Y(\nu,a)\defeq y_{40}(1;\nu;a).
\]
\end{lemma}
\begin{proof}
Consider a connected graph with $L$ loops, $I$ lines and vertices $v$ of lattice degree $d_v$ ($1$ for the current vertex, $0$ for the seagull and the quartic vertex, $2-n_v$ for the $n_v$-photon vertex of the hopping term), at zero external momentum. Substitute $q=tq'$ in every loop momentum: $q\in\BZ_\eps$ iff $q'\in\BZ_{t\eps}$, $d^3q=t^3d^3q'$, and $\frac4{\eps^2}\sin^2\frac{\eps tq'_\mu}2=t^2\cdot\frac4{(t\eps)^2}\sin^2\frac{(t\eps)q'_\mu}2$, so $\hat q^2$ at spacing $\eps$ equals $t^2$ times $\hat{q'}^2$ at spacing $t\eps$. Each propagator therefore gives $t^{-2}$: $(\hat q^2+s)^{-1}=t^{-2}(\hat{q'}^2+s/t^2)^{-1}$ and $(\hat k^2+\lmu^2)^{-1}=t^{-2}(\hat{k'}^2+\lmu^2/t^2)^{-1}$. A vertex of the hopping term is $\eps^{n_v-2}w(\eps q)$ with $w$ a trigonometric polynomial, and $\eps^{n_v-2}w(\eps q)=t^{2-n_v}(t\eps)^{n_v-2}w\big((t\eps)q'\big)$, i.e.\ $t^{d_v}$ times the same vertex at spacing $t\eps$; the quartic vertex is a constant. Hence the amplitude obeys $A(s,\lmu,\eps)=t^{3L-2I+\sum_vd_v}A(s/t^2,\lmu/t,t\eps)$, and for a two-point graph of order $n$ the exponent is the superficial degree $3-\frac E2-n=2-n$ ($E=2$). A counterterm vertex $d_{\alpha'\beta'}$ of lower order $n'$ inserted on a line contributes, by induction on the order, the factor $t^{2-n'}$ and one more propagator $t^{-2}$, which leaves the exponent $2-n$ unchanged. $d_{\alpha\beta}$ is the zero-momentum value of such a sum \cite[(3.32)--(3.33)]{K1}, so the identity holds. Differentiate in $s$ at fixed $t$ and put $t=\sqrt s$: $y_{\alpha\beta}(s;\lmu;\eps)=s^{-n/2}y_{\alpha\beta}(1;\lmu/\sqrt s;\eps\sqrt s)$, and $n=2$ at $(4,0)$.
\end{proof}

\paragraph{Conventions at unit scalar mass.} For the rest of the proof the spacing is $a\in(0,1]$, $\BZ_a\defeq[-\pi/a,\pi/a)^3$, $\langle f\rangle_a\defeq\int_{\BZ_a}f\,\frac{d^3k}{(2\pi)^3}$, $\langle\langle\cdot\rangle\rangle_a$ for two loop momenta; $\hat e(k)\defeq\sum_\mu\frac4{a^2}\sin^2\frac{ak_\mu}2$, so that $\frac4{\pi^2}|k|^2\le\hat e(k)\le|k|^2$ on $\BZ_a$ (\cref{lem:a05:el}(c)) and $\hat e\le12/a^2$; $\chi(k)\defeq\hat e(k)+1$, $C(q)\defeq\chi(q)^{-1}$ (scalar mass $1$), $D_\nu(k)\defeq(\hat e(k)+\nu^2)^{-1}$ (photon mass $\nu>0$; the propagator is diagonal in the photon indices), $D_0\defeq1/\hat e$. In the per-component normalisation of (i)--(v) the current vertex carries $\tilde V_\mu(p,p')\defeq-\frac2a\sin\frac{a(p+p')_\mu}2$ and the phase $e^{-\mathrm iak_\mu/2}$, $k=p'-p$ the incoming photon momentum; the seagull is $e^{-\mathrm iak_\mu/2}\cos\frac{a(p+p')_\mu}2\,\delta_{\mu\nu}$; the three- and four-photon vertices $V_3$, $V_4$ of the hopping term are $a$ and $a^2$ times trigonometric polynomials of $a(p+p')$ and $ak$ (\cref{lem:a05:deriv}). Numerical factors of the Feynman rules play no role and are suppressed. Soft and hard momenta: $Q\defeq[-\frac12,\frac12]^3\subset\BZ_a$, $Q_s\defeq\{|k|\le\frac12\}$, hard $\defeq\BZ_a\setminus Q_s$. For hard $k$, $\hat e(k)\ge\pi^{-2}$, so $D_\nu\le(1+\pi^2)D_1$; for every $k$, $D_\nu\le D_0\le\frac{\pi^2}4|k|^{-2}$ on $\BZ_a$. In particular, for every $\nu>0$ and $a\le1$,
\begin{equation}\label{eq:a05:soft}
\sigma_1\defeq\langle\mathbf 1_{Q_s}D_\nu\rangle_a\le\tfrac1{16},\qquad\langle\mathbf 1_{Q_s}D_\nu^2|k|^2\rangle_a\le\tfrac{\pi^2}{64},\qquad\langle C^2\rangle_a\le\tfrac{\pi^2}{64},\qquad\langle\chi^{-3/2}\rangle_a\le C\big(1+\ln\tfrac1a\big),
\end{equation}
by $\frac{\pi^2}4\int_{|k|\le1/2}\frac{d^3k}{(2\pi)^3|k|^2}=\frac1{16}$, $(\frac{\pi^2}4)^2\int_{|k|\le1/2}\frac{d^3k}{(2\pi)^3|k|^2}=\frac{\pi^2}{64}$, $(\frac{\pi^2}4)^2\int_{\R^3}\frac{d^3k}{(2\pi)^3(|k|^2+\pi^2/4)^2}=\frac{\pi^2}{64}$ and $(\frac{\pi^2}4)^{3/2}\int_{|k|\le\sqrt3\pi/a}\frac{d^3k}{(2\pi)^3(|k|^2+\pi^2/4)^{3/2}}$. Throughout, $C$ denotes a constant depending only on $a_0$ (fixed in \cref{lem:a05:H}), never on $\nu$, $a$ or $s$; it may change from line to line.

\paragraph{Subtraction operators, degrees, chords.} For a one-particle-irreducible subgraph $\gamma$ of $\Gamma$ with external momenta $q^\gamma$ and internal momenta $k^\gamma$, the subtraction operator of order $\delta\in\{0,1\}$ acts on the integrand $F_\gamma(k^\gamma;q^\gamma)$ at fixed $k^\gamma$ \cite[Def.~3.1]{Rei88b}:
\[
t^0_\gamma F_\gamma=F_\gamma(k^\gamma;0),\qquad t^1_\gamma F_\gamma=F_\gamma(k^\gamma;0)+\sum_\rho P_\rho(q^\gamma)\,\partial_{q_\rho}F_\gamma(k^\gamma;0),\qquad P_\rho(q)=a^{-1}p_\rho(aq),
\]
with $p_\rho$ smooth and $2\pi$-periodic, $p_\rho(x)=x_\rho+O(|x|^2)$ (\cite[Def.~3.1, (3-2), and App.~A]{Rei88b}; the same $P_\rho$ is used throughout). These operators preserve periodicity and act as the Taylor operators do on the degrees. The degree of $\gamma$ is $\omega(\gamma)=3L(\gamma)-2I(\gamma)+\sum_{v\in\gamma}d_v=3-\frac{E(\gamma)}2-n(\gamma)$ ($E$ external legs, $n$ the order); it is the UV-divergence degree of \cite[(2-28)--(2-30)]{Rei88b} in three dimensions, computed with the lattice degrees of \cite[Defs.~2.2--2.3]{Rei88}, the lattice degree of every vertex function being its continuum dimension. The lines of $\Gamma$ are ordered with the photon lines first, and the chord set of every connected subgraph is the one determined by this order through the rule of \cite[(2-24)]{Rei88b}: the chords of $\Gamma$ carry the loop momenta, all other line momenta follow by conservation, the chords of a subgraph $\gamma$ are chords of $\Gamma$, and the internal momenta of $\gamma$ are its chord momenta, which carry no external momentum of $\gamma$ \cite[(2-16)]{Rei88b}.

\begin{lemma}[the forest form of King's coefficient]\label{lem:a05:forest}
On $\eps\Z^3$, at any spacing and masses, with signed combinatorial weights $c_\Gamma$ whose values play no role below,
\[
d_{40}=\sum_\Gamma c_\Gamma\,\langle\langle R'F_\Gamma\rangle\rangle,\qquad y_{40}=\partial_sd_{40}=-\sum_\Gamma\sum_{l}c_\Gamma\,\langle\langle R'F_{\Gamma,\mathrm{dbl}(l)}\rangle\rangle ,
\]
where $\Gamma$ runs over the twelve one-particle-irreducible two-point graphs of order $e^4$ at zero external momentum (listed next), $F_\Gamma$ is the integrand of the Feynman rules without counterterms, $R'\defeq1-\sum_\gamma t^{\omega(\gamma)}_\gamma$ with $\gamma$ running over the proper one-particle-irreducible subgraphs with $\omega(\gamma)\ge0$, $l$ runs over the scalar lines of $\Gamma$, and $F_{\Gamma,\mathrm{dbl}(l)}$ is $F_\Gamma$ with the propagator of $l$ squared.
\end{lemma}
\begin{proof}
King's scheme \cite[(3.32)--(3.33)]{K1} has scalar mass counterterms at the orders $2\le\alpha+2\beta\le4$, each fixed by the vanishing at zero momentum of the one-particle-irreducible two-point function of its order with the lower-order counterterms inserted, and no photon or vertex counterterm. Compare it with the scheme in which every one-particle-irreducible type of non-negative degree receives the counterterm $-t^{\omega}$ of its zero-momentum Taylor polynomial, computed recursively in the order; its amplitudes are the zone integrals of the forest-subtracted integrands \cite[Sect.~4]{Rei88b}, and since the forests of proper divergent subgraphs have at most one element here (property (iv) below), the forest operator is $R'$. The proper divergent subgraphs of the graphs of order $e^4$ are of three types: (a) the one-loop scalar self-energy of order $e^2$ (exchange graph, seagull tadpole; $\omega=1$); (c) the one-loop vacuum polarisation of order $e^2$ (bubble, seagull loop; $\omega=1$); (d) the one-loop $\phi\bar\phi A$ vertex of order $e^3$ (triangle, two swordfish graphs, $A^3$-tadpole; $\omega=0$). The two schemes agree on $\eps\Z^3$:
(a) the order-$1$ part $t^1_\gamma-t^0_\gamma$ of the exchange integrand $f(q;p)=\widehat{(p+q)}^2D(p-q)C(q)$ is $\sum_\rho P_\rho(p)\partial_{p_\rho}f(q;0)$, which is odd in the loop momentum $q$ ($D$, $C$ even), while the seagull-tadpole integrand $D(k)\cos(ap_\mu)$ has zero gradient at $p=0$; so the order-$1$ term integrates to zero identically in the remaining momenta, and the order-$0$ term is King's $d_{20}$-insertion.
(c) The one-loop polarisation $\Pi^{(20)}_{\mu\nu}(k)$ (bubble plus seagull loop) satisfies $\Pi^{(20)}_{\mu\nu}(0)=\pm\big\langle\partial_\nu\big[\tfrac2a\sin(aq_\mu)\,C(q)\big]\big\rangle_a=0$ for all $\mu,\nu$, the integral of a derivative of a periodic function (on the diagonal the seagull loop is the derivative of the current loop; off the diagonal the bubble alone is a derivative; this is the lattice Ward identity used after \eqref{eq:a05:Pi}), and $\nabla_k\Pi^{(20)}(0)=0$: writing $\Pi^{(20)}_{\mu\nu}(k)=e^{-\mathrm ia(k_\mu-k_\nu)/2}T_{\mu\nu}(k)$, $T$ is even in $k$ (substitute $q\mapsto-q$), so $\nabla T(0)=0$ and $\nabla_k\Pi^{(20)}(0)=-\frac{\mathrm ia}2(e_\mu-e_\nu)\Pi^{(20)}_{\mu\nu}(0)=0$. Hence $t^1_k\Pi^{(20)}=0$: the photon counterterm of the forest scheme vanishes, as King's does.
(d) The vertex counterterm is the value at zero momenta of the sum of the one-loop vertex graphs, and each vanishes by $q\mapsto-q$: in the triangle $\tilde V_\mu(q,q)=-\frac2a\sin(aq_\mu)$ is odd, the other factors are even and the two phases of the internal photon cancel; the two swordfish graphs sum to $\langle2\mathrm i\sin(aq_\mu)D(q)C(q)\rangle_a=0$; the $A^3$-tadpole is proportional to $\sin\frac{a(p+p')_\mu}2=0$ at $p=p'=0$. So neither scheme has a vertex counterterm.
Therefore the zero-momentum value of the order-$e^4$ two-point function with King's insertions is $\sum_\Gamma c_\Gamma\langle\langle R'F_\Gamma\rangle\rangle$. Differentiating in $s$: the operators $t_\gamma$ act on momenta at fixed masses, so $\partial_s$ commutes with $R'$, and $\partial_sC=-C^2$ doubles one scalar line at a time.
\end{proof}

\paragraph{The twelve graphs.} They are the one-particle-irreducible two-point placements of the vertex sets $V_1^4$, $V_2V_1^2$, $V_2^2$, $V_3V_1$, $V_4$; every other placement contains a bridge. With the names of the table in \cref{lem:a05:40}:
\begin{center}\small
\begin{tabular}{@{}cllll@{}}
& vertices & graph & proper subgraphs with $\omega\ge0$ (type; photon lines inside) & class; removed terms of \eqref{eq:a05:OS}\\\hline
$\Gamma_0$ & $V_1^4$ & rainbow & inner exchange (a; 1) & A; $l=$ inner line\\
$\Gamma_1$ & $V_1^4$ & exchange, bubble inserted & bubble (c; 0) & B; $l$ in the bubble\\
$\Gamma_2$ & $V_1^4$ & crossed & two triangles (d; 1), overlapping & A; coefficient $0$\\
$\Gamma_3$ & $V_2V_1^2$ & exchange, seagull tadpole on the inner line & seagull tadpole (a; 1) & A; none\\
$\Gamma_4$ & $V_2V_1^2$ & seagull tadpole, bubble inserted & bubble (c; 0) & B; $l$ in the bubble\\
$\Gamma_5$ & $V_2V_1^2$ & seagull at an end (two mirror graphs) & swordfish vertex (d; 1) & A; coefficient $0$\\
$\Gamma_6$ & $V_2V_1^2$ & exchange, seagull loop inserted & seagull loop (c; 0) & B; vanishes, see (0)\\
$\Gamma_7$ & $V_2V_1^2$ & seagull in the middle & two swordfish vertices (d; 1), sharing $V_2$ & A; coefficient $0$\\
$\Gamma_8$ & $V_3V_1$ & $V_3$-tadpole vertex (two orientations) & $A^3$-tadpole (d; 1), no scalar line & A; none\\
$\Gamma_9$ & $V_2^2$ & seagull tadpole, seagull loop inserted & seagull loop (c; 0) & B; vanishes, see (0)\\
$\Gamma_{10}$ & $V_2^2$ & melon & none & A; none\\
$\Gamma_{11}$ & $V_4$ & $V_4$ double tadpole & none & no scalar line
\end{tabular}
\end{center}
By inspection of the twelve graphs: (i) $\omega(\Gamma)=0$, and for every scalar line $l$, $\omega_{\mathrm{dbl}}(\Gamma)=-2$, where $\omega_{\mathrm{dbl}}(\gamma)\defeq\omega(\gamma)-2\cdot\mathbf 1[l\in\gamma]$ is the degree of the doubled graph (the doubled line is a two-point insertion of degree $0$); (ii) every proper one-particle-irreducible subgraph with $\omega\ge0$ is one-loop, of type (a), (c) or (d), and contains at most one photon line, which, the photon lines being first in the order, is that subgraph's own chord; the subgraphs of type (c) contain no photon line; (iii) no subgraph containing $l$ has $\omega_{\mathrm{dbl}}\ge0$; (iv) two proper divergent subgraphs of the same $\Gamma$ always share a line or a vertex (the triangles of $\Gamma_2$ overlap, the swordfish vertices of $\Gamma_7$ share the seagull), so the forests of proper subgraphs in the sense of \cite[Sect.~2]{Rei88b} (members nested or with no common lines and vertices) have at most one element; (v) in class A both photon lines are chords of $\Gamma$, with independent loop momenta $k_1,k_2$, and the scalar lines carry integer combinations of them; in class B one photon line is a chord with momentum $k$, the other photon line carries $\pm k$, and the second chord is a scalar line of the closed scalar loop, with momentum $q$.

By (iv), $R'=1-\sum_\gamma t_\gamma$ splits according to whether $\gamma$ contains the doubled line: for every $\Gamma$ and scalar line $l$, pointwise,
\begin{equation}\label{eq:a05:OS}
R'F_{\Gamma,\mathrm{dbl}(l)}=R_{\mathrm{dbl}}F_{\Gamma,\mathrm{dbl}(l)}-\sum_{\gamma\ni l,\ \omega(\gamma)\ge0}\big(t^{\omega(\gamma)}_\gamma F_{\gamma,\mathrm{dbl}(l)}\big)\,F_{\bar\Gamma_\gamma},
\qquad R_{\mathrm{dbl}}\defeq1-\sum_{\gamma\not\ni l,\ \omega(\gamma)\ge0}t^{\omega(\gamma)}_\gamma ,
\end{equation}
where $F_{\bar\Gamma_\gamma}$ is the integrand of the lines and vertices of $\Gamma$ outside $\gamma$. $R_{\mathrm{dbl}}$ is the forest formula of the doubled graph with the labels $\omega_{\mathrm{dbl}}$ on every one-particle-irreducible subgraph: by (iii) the subgraphs containing $l$, and $\Gamma$ itself, have negative label and are not subtracted. The subtracted terms (``removed terms'') are polynomials of degree $\le\omega(\gamma)$ in $P(q^\gamma)$ whose coefficients are integrals over $k^\gamma$.

\begin{lemma}[photon factorisation]\label{lem:a05:PF}
Fix $\Gamma$, a scalar line $l$ and the photon mass $\nu$, and write $F^{(\nu)}$ for the integrand at that photon mass. In class A, $R_{\mathrm{dbl}}F^{(\nu)}_{\Gamma,\mathrm{dbl}(l)}(k_1,k_2)=D_\nu(k_1)D_\nu(k_2)\,R_{\Gamma,l}(k_1,k_2)$, and each removed term of \eqref{eq:a05:OS} is $D_\nu(k^\gamma)D_\nu(k^{\bar\Gamma})$ times a $\nu$-free function. In class B, $R_{\mathrm{dbl}}F^{(\nu)}_{\Gamma,\mathrm{dbl}(l)}(k,q)=D_\nu(k)^2R_{\Gamma,l}(k,q)$, and the same holds for $R'F^{(\nu)}_{\mathrm{dbl}}$ and for the removed term. The functions $R_{\Gamma,l}$ and the cofactors do not depend on $\nu$.
\end{lemma}
\begin{proof}
Every factor of $F_\Gamma$ other than the two photon propagators is $\nu$-free: the vertex functions are mass-free trigonometric polynomials and the scalar lines have mass $1$. A subtraction operator $t_\gamma$ Taylor-expands the integrand of $\gamma$ in $\gamma$'s external momenta at fixed internal momenta, and $\gamma$'s internal momentum is its chord momentum, which carries no external momentum of $\gamma$ \cite[(2-16), (2-24)]{Rei88b}. By (ii) every subtracted $\gamma$ contains either exactly one photon line, which is $\gamma$'s own chord, so that its propagator is a function of $\gamma$'s internal momentum alone and is untouched by $t_\gamma$, or no photon line (type (c)), its photon legs being external to $\gamma$, so that the attached photon propagators belong to $\bar\Gamma_\gamma$ and are cofactors. Hence no subtraction operator acts on a photon propagator, and the photon propagators factor out of every term of \eqref{eq:a05:OS} with the momenta of (v).
\end{proof}

\begin{lemma}[the fixed-mass bound]\label{lem:a05:H}
There are $a_0\in(0,1]$ and $K<\infty$ such that, at scalar and photon mass $1$, for every $\Gamma$ and every scalar line $l$ of $\Gamma$,
\[
\sum_{|\beta|\le6}\ \sup_{0<a\le a_0}\ \big\langle\big\langle\big|\partial^\beta R_{\mathrm{dbl}}F^{(1)}_{\Gamma,\mathrm{dbl}(l)}\big|\big\rangle\big\rangle_a\le K,
\]
$\partial^\beta$ running over the derivatives in the two loop momenta.
\end{lemma}
\begin{proof}
$R_{\mathrm{dbl}}F^{(1)}$ is a forest-subtracted lattice Feynman integrand in the sense of \cite{Rei88b}: the propagators are of the form $(\hat e(l)+\mu^2)^{-1}$ with $\mu=1$, the numerators are trigonometric polynomials of the class of \cite[Defs.~2.1--2.5]{Rei88} whose lattice degrees equal their continuum dimensions, the doubled line is a two-point insertion of degree $0$, and the labels $\delta(\gamma)=\omega_{\mathrm{dbl}}(\gamma)$ are the degrees of the doubled graph itself, so the conditions (4-3) and (4-4) of \cite{Rei88b} hold (remark after (4-4) there), with $t^\delta=0$ for $\delta<0$. The chord assignment obeys \cite[(2-24)]{Rei88b}, and the line momenta are natural in the sense of \cite[Def.~3.1]{Rei88}: every subtracted subgraph is one-loop, so its internal momentum is, up to sign, the momentum of one of its lines, and the set of line momenta of a connected graph without cut lines is natural (\cite{Rei88}, remark after Def.~3.1). By the renormalization theorem of \cite[Sects.~4--5]{Rei88b}, whose proof applies verbatim in three dimensions (the dimension enters only through the degree count; cf.\ the remark before Def.~2.5 of \cite{Rei88}), the integrand $R_{\mathrm{dbl}}F^{(1)}$ is $\frac{2\pi}a$-periodic, lies in the class of \cite{Rei88} and has negative degree on every Zimmermann subspace, the overall degree being $\omega_{\mathrm{dbl}}(\Gamma)=-2$; by \cite[Lemmas~2.1--2.2]{Rei88} the derivatives $\partial^\beta$ stay in the class and do not raise any of these degrees. The three-dimensional form of the renormalization theorem is read off its proof as follows. The dimension $d$ enters that proof at three places only: through the term $d\,m(\gamma)$, $m(\gamma)$ the number of loops, in the divergence degree \cite[(2-28), (2-31)]{Rei88b}; through the quantities $M_U(\gamma)=d\sum m$ of \cite[(5-10)]{Rei88b} and the bound $M_U(\Gamma)\ge dh$ of \cite[App.~C]{Rei88b}, which is $d$ times an inequality between loop counts and ranks of the coefficient matrices of the loop-momentum parametrisation of the Zimmermann subspace, the same for every $d$; and through the target inequality \cite[(5-5)]{Rei88b}, $dh+\operatorname{degr}_tR_\Gamma<0$. The degree bounds (5-11)--(5-13) of \cite[Lemma~5.2]{Rei88b} are proved by induction over the forest from the lattice degrees of the single propagators and vertex functions, which do not involve the dimension, from the subtraction degrees and from the loop counts, which add under the reduction of a subgraph, the same way as the degrees of (2-28) do; so with the degrees counted in $d$ dimensions and labels satisfying (4-3) and (4-4) in those degrees, every line of the induction holds for every $d$, and its conclusion in $d=3$, $3h+\operatorname{degr}_tR_\Gamma<0$ on every Zimmermann subspace, is the hypothesis (3.9) of the power counting theorem of \cite{Rei88}. Everything else in the proof (the subtraction operators and their degree calculus, the forest algebra, the periodicity of $R_\Gamma$) treats the momenta as lists of real variables and does not depend on the dimension. The power counting theorem \cite[Thm.~1]{Rei88} applies to each $\partial^\beta R_{\mathrm{dbl}}F^{(1)}$, and its proof bounds the integrand on every sector by non-negative majorants whose integrals over the zone are bounded uniformly in $a\le a_0$ when all degrees are negative (\cite[Sect.~4 and Thm.~2]{Rei88}), which is the uniform absolute bound claimed; $a_0$ and $K$ depend on the masses, here $(1,1)$, and on the finitely many integrands $(\Gamma,l,\beta)$.
\end{proof}

% a05_twoloop_v2: proof of lem:a05:D45 (v), part 2: derivative bounds, box inequality, polarisation remainder, the unit-mass bound, conclusion (input from a05_cnum_NEG.tex)
\begin{lemma}[derivative bounds at unit mass]\label{lem:a05:EB}
Let $a\le1$ and $q,k\in\R^3$.
\begin{enumerate}[label=\textup{(\roman*)},nosep,leftmargin=*]
\item $\big|\sqrt{\hat e(q+k)}-\sqrt{\hat e(q)}\big|\le|k|$; hence $C(q+k)\le2C(q)$ and $\chi(q+k)\le2\chi(q)$ for $|k|\le\frac12$.
\item $a^2\le13/\chi(q)$.
\item $|\partial_\mu\hat e|\le2\sqrt{\hat e}$, and every derivative of $\hat e$ of order $\ge2$ is bounded by $2$. Consequently, for $|k|\le\frac12$ and all multi-indices $b,\beta$: $|\partial^b_q\partial^\beta_kC(q+k)|\le C_{b,\beta}\,\chi(q)^{-1-(|b|+|\beta|)/2}$, $|\partial^b_q\partial^\beta_k\tilde V_\mu(q,q+k)|\le C_{b,\beta}\,\chi(q)^{(1-|b|-|\beta|)/2}$, and the derivatives of $\cos\frac{a(2q+k)_\mu}2$ and of the phases $e^{\pm\mathrm iak_\mu/2}$ are bounded by $C_{b,\beta}\,\chi(q)^{-(|b|+|\beta|)/2}$: each derivative gains a factor $\chi(q)^{-1/2}$ relative to the bound of the undifferentiated factor ($C(q+k)\le2\chi(q)^{-1}$, $|\tilde V_\mu(q,q+k)|\le2\sqrt{\hat e(q)}+|k|\le3\chi(q)^{1/2}$, $|\cos|,|e^{\mathrm i\cdot}|\le1$).
\item $|P_\rho(k)|\le C\sqrt{\hat e(k)}$ for all $k$, and $|P_\rho(k)-k_\rho|\le Ca|k|^2$ for $|k|\le\frac12$.
\end{enumerate}
\end{lemma}
\begin{proof}
(i) With $v_\mu(q)\defeq\frac2a\sin\frac{aq_\mu}2$, $\sqrt{\hat e}=|v|$ and $|v_\mu(q+k)-v_\mu(q)|\le|k_\mu|$ by $|\sin x-\sin y|\le|x-y|$; the triangle inequality gives the first claim. If $x\defeq\sqrt{\hat e(q)}\ge\frac12$, then $\hat e(q+k)+1\ge(x-\frac12)^2+1\ge\frac{x^2+1}2$ (equivalent to $(x-1)^2+\frac12\ge0$); if $x<\frac12$, then $\hat e(q+k)+1\ge1\ge\frac{x^2+1}2$. Also $\hat e(q+k)+1\le(x+\frac12)^2+1\le2x^2+2$.
(ii) $\hat e\le12/a^2$ and $a\le1$.
(iii) $\partial_\mu\hat e=\frac2a\sin(aq_\mu)=\frac4a\sin\frac{aq_\mu}2\cos\frac{aq_\mu}2$, of modulus $\le2|v_\mu|$; $\partial_\mu\partial_\nu\hat e=2\cos(aq_\mu)\delta_{\mu\nu}$, and each further derivative multiplies by $a\le1$. A derivative of $C=\chi^{-1}$ produces $C^2\partial\hat e$ with $|\partial\hat e|\le2\chi^{1/2}$, a gain of $2\chi^{-1/2}$; a derivative of a factor $\partial\hat e$ gives a factor $\le2\le2\chi^{1/2}\chi^{-1/2}$. The $j$-th derivative of $\tilde V_\mu(q,q+k)=-\frac2a\sin\frac{a(2q+k)_\mu}2$ is bounded by $2a^{j-1}\le2\cdot13^{(j-1)/2}\chi(q)^{-(j-1)/2}$ for $j\ge1$ by (ii), and that of $\cos\frac{a(2q+k)_\mu}2$ or $e^{\pm\mathrm iak_\mu/2}$ by $a^j\le13^{j/2}\chi(q)^{-j/2}$. For the factor $C(q+k)$ and its derivatives use $\chi(q+k)\ge\frac12\chi(q)$ from (i). Induction on $|b|+|\beta|$ gives the stated bounds.
(iv) Applying \cite[(3-2)]{Rei88b} to $F\equiv q_\rho$ gives $P_\rho(0)=0$ and $\nabla P_\rho(0)=e_\rho$, so $p_\rho(x)=x_\rho+O(|x|^2)$, and by periodicity $p_\rho$ vanishes on $2\pi\Z^3$. Hence $|p_\rho(x)|\le\|\nabla p_\rho\|_\infty\dist(x,2\pi\Z^3)\le\|\nabla p_\rho\|_\infty\frac\pi2\big(4\sum_\mu\sin^2\frac{x_\mu}2\big)^{1/2}$ by \cref{lem:a05:el}(c), and $|P_\rho(k)|=a^{-1}|p_\rho(ak)|\le C\sqrt{\hat e(k)}$; and $|P_\rho(k)-k_\rho|=a^{-1}|p_\rho(ak)-ak_\rho|\le a^{-1}C|ak|^2=Ca|k|^2$ for $|ak|\le\frac12$.
\end{proof}

\begin{lemma}[box inequality]\label{lem:a05:box}
\textup{(i)} For $g\in C^d$ on a cube $Q_d\subset\R^d$ of unit side, $\sup_{Q_d}|g|\le\sum_{S\subset\{1,\dots,d\}}\int_{Q_d}|\partial_Sg|$, $\partial_S\defeq\prod_{i\in S}\partial_i$ ($\partial_\emptyset g=g$). \textup{(ii)} For $G$ smooth and $\frac{2\pi}a$-periodic in $(k_1,k_2)\in\R^6$,
\[
\sup_{k_2\in Q}\big\langle|G(\cdot,k_2)|\big\rangle_{a,k_1}\le\sum_{S\subset\{1,2,3\}}\int_Q\big\langle|\partial^{k_2}_SG(\cdot,k_2)|\big\rangle_{a,k_1}\,dk_2,\qquad
\sup_{Q^2}|G|\le\sum_{S\subset\{1,\dots,6\}}\int_{Q^2}|\partial_SG| .
\]
\end{lemma}
\begin{proof}
(i) For $x,y\in Q_d$ the fundamental theorem of calculus, applied coordinate by coordinate, gives $g(x)=\sum_S\int_{[y_S,x_S]}(\partial_Sg)(t_S,y_{S^c})\,dt_S$, so $|g(x)|\le\sum_S\int_{[0,1]^{|S|}}|(\partial_Sg)(t_S,y_{S^c})|\,dt_S$ for every $y$; average over $y\in Q_d$. (ii) Apply (i) in $k_2$ at fixed $k_1$ and integrate over $k_1$ (Fubini), respectively apply (i) in the six variables.
\end{proof}
Applied to $G=\chi(k_2)R_{\mathrm{dbl}}F^{(1)}$, respectively $G=\chi(k_1)\chi(k_2)R_{\mathrm{dbl}}F^{(1)}$: on $Q$, $\chi\le\frac74$, $|\partial_i\chi|=|\frac2a\sin(ak_i)|\le2|k_i|\le1$, and the mixed derivatives of $\chi$ vanish ($\hat e$ is a sum of functions of one coordinate each), so by the Leibniz rule both right sides of (ii) are $\le CK$, with $K$ of \cref{lem:a05:H} (derivatives of order $\le3$, respectively $\le6$) and a numerical $C$.

\begin{lemma}[the polarisation remainder]\label{lem:a05:VPR}
Let $F_b(q;k)\defeq e^{-\mathrm ia(k_\mu-k_\nu)/2}\tilde V_\mu(q,q+k)\tilde V_\nu(q+k,q)C(q)C(q+k)$ be the bubble integrand at unit mass with photon indices $\mu,\nu$, or either of its doubled versions ($C(q)\to C(q)^2$ or $C(q+k)\to C(q+k)^2$). For $a\le1$ and $|k|\le\frac12$,
\[
\big|\big\langle(1-t^1_k)F_b(\cdot;k)\big\rangle_a\big|\le C_{\mathrm{VP}}|k|^2,
\]
and the left side vanishes at $k=0$. Moreover $\Pi_{\mathrm{dbl}}(0)\defeq\langle F_{b,\mathrm{dbl}}(\cdot;0)\rangle_a$ obeys $|\Pi_{\mathrm{dbl}}(0)|\le4\langle C^2\rangle_a\le\frac{\pi^2}{16}$ and $\nabla_k\langle F_{b,\mathrm{dbl}}(\cdot;k)\rangle_a\big|_{k=0}=-\frac{\mathrm ia}2(e_\mu-e_\nu)\Pi_{\mathrm{dbl}}(0)$, of modulus $\le a|\Pi_{\mathrm{dbl}}(0)|$.
\end{lemma}
\begin{proof}
For fixed $q$, $(1-t^1_k)F_b(q;k)=\big[F_b(q;k)-F_b(q;0)-k\cdot\nabla_kF_b(q;0)\big]+\big(k-P(k)\big)\cdot\nabla_kF_b(q;0)$. The bracket equals $\int_0^1(1-\theta)\,k^{\mathsf T}\nabla_k^2F_b(q;\theta k)\,k\,d\theta$, and by \cref{lem:a05:EB}(iii), $|\partial^\beta_kF_b(q;k')|\le C\chi(q)^{-1-|\beta|/2}$ for $|k'|\le\frac12$ (two vertex factors of size $\chi^{1/2}$, two propagators $\chi^{-1}$, each derivative gaining $\chi^{-1/2}$; the doubled versions are smaller by a factor $2\chi^{-1}$). Hence $\langle|\nabla^2_kF_b(\cdot;\theta k)|\rangle_a\le C\langle\chi^{-2}\rangle_a\le C$ by \eqref{eq:a05:soft}, and the bracket contributes $\le C|k|^2$. For the second term, $|k-P(k)|\le Ca|k|^2$ by \cref{lem:a05:EB}(iv) and $\langle|\nabla_kF_b(\cdot;0)|\rangle_a\le C\langle\chi^{-3/2}\rangle_a\le C(1+\ln\frac1a)$ by \eqref{eq:a05:soft}; and $a(1+\ln\frac1a)\le1$ for $a\le1$. At $k=0$ the remainder vanishes identically in $q$, since $P(0)=0$. For the last claims: at $k=0$ the integrand is $\tilde V_\mu(q,q)\tilde V_\nu(q,q)C(q)^3$ with $|\tilde V_\mu(q,q)|=\frac2a|\sin(aq_\mu)|\le2\sqrt{\hat e(q)}$, so $|\Pi_{\mathrm{dbl}}(0)|\le4\langle\hat eC^3\rangle_a\le4\langle C^2\rangle_a$; and $\langle F_{b,\mathrm{dbl}}(\cdot;k)\rangle_a=e^{-\mathrm ia(k_\mu-k_\nu)/2}T_{\mathrm{dbl}}(k)$ with $T_{\mathrm{dbl}}$ even in $k$ (the substitution $q\mapsto-q$ maps the integrand at $k$ to the integrand at $-k$, the doubled line to itself), so $\nabla T_{\mathrm{dbl}}(0)=0$.
\end{proof}

\begin{proposition}[the unit-mass bound, uniform in the photon mass]\label{prop:a05:U}
With $a_0$ of \cref{lem:a05:H} there is $C_U<\infty$, depending on $a_0$ only, such that $|Y(\nu,a)|\le C_U$ for all $\nu>0$ and $a\in(0,a_0]$.
\end{proposition}
\begin{proof}
By \cref{lem:a05:forest} at $s=1$, $Y(\nu,a)$ is a finite sum of terms $c_\Gamma\langle\langle R'F^{(\nu)}_{\Gamma,\mathrm{dbl}(l)}\rangle\rangle_a$. Each is bounded by a constant of the stated kind.

\emph{(0) Vanishing terms.} $\Gamma_{11}$ has no scalar line. For $\Gamma_6$ and $\Gamma_9$ the polarisation subgraph is the seagull loop; the seagull vertex depends on its two photon momenta only through their sum, which is $0$ at zero external momentum, so the loop integrand $F_{\mathrm{VP}}(q;k)$ does not depend on $k$ and $(1-t^1_k)F_{\mathrm{VP}}\equiv0$. Hence $R'F_{\mathrm{dbl}}\equiv0$ for every $l$: if $l$ lies on the path, $R'=R_{\mathrm{dbl}}$ subtracts the loop; if $l$ lies in the loop, $R'F_{\mathrm{dbl}}=(1-t^1_k)F_{\mathrm{dbl}}=0$.

\emph{(1) Class A ($\Gamma_0,\Gamma_2,\Gamma_3,\Gamma_5,\Gamma_7,\Gamma_8,\Gamma_{10}$), the $R_{\mathrm{dbl}}$ part.} By \cref{lem:a05:PF}, $|R_{\mathrm{dbl}}F^{(\nu)}|=D_\nu(k_1)D_\nu(k_2)|R_{\Gamma,l}(k_1,k_2)|=D_\nu(k_1)\chi(k_1)\,D_\nu(k_2)\chi(k_2)\,|R_{\mathrm{dbl}}F^{(1)}|$. Split $\BZ_a^2$ according to whether $k_1$, $k_2$ are hard or soft. On hard momenta $D_\nu\chi\le1+\pi^2$. Hence
\begin{align*}
\langle\langle\mathbf 1_{hh}|R_{\mathrm{dbl}}F^{(\nu)}|\rangle\rangle_a&\le(1+\pi^2)^2\langle\langle|R_{\mathrm{dbl}}F^{(1)}|\rangle\rangle_a\le(1+\pi^2)^2K,\\
\langle\langle\mathbf 1_{hs}|R_{\mathrm{dbl}}F^{(\nu)}|\rangle\rangle_a&\le(1+\pi^2)\,\sigma_1\sup_{k_2\in Q}\big\langle|\chi(k_2)R_{\mathrm{dbl}}F^{(1)}(\cdot,k_2)|\big\rangle_{a,k_1}\le(1+\pi^2)\sigma_1CK,
\end{align*}
by \cref{lem:a05:box}(ii) and \cref{lem:a05:H}, and the same for $sh$ with $k_1\leftrightarrow k_2$; and $\langle\langle\mathbf 1_{ss}|R_{\mathrm{dbl}}F^{(\nu)}|\rangle\rangle_a\le\sigma_1^2\sup_{Q^2}|\chi(k_1)\chi(k_2)R_{\mathrm{dbl}}F^{(1)}|\le\sigma_1^2CK$. With $\sigma_1\le\frac1{16}$ from \eqref{eq:a05:soft}, $|\langle\langle R_{\mathrm{dbl}}F^{(\nu)}\rangle\rangle_a|\le C$.

\emph{(2) Class A, the removed terms of \eqref{eq:a05:OS}.} Type (d) ($\Gamma_2$, $\Gamma_5$, $\Gamma_7$): the coefficient $\langle F_{\gamma,\mathrm{dbl}}(\cdot;0)\rangle_a$ is the one-loop vertex subgraph with the line $l$ doubled at zero external momenta, and it vanishes by $q\mapsto-q$: for the triangle the integrand is $\sum_\alpha\tilde V_\alpha(0,q)\tilde V_\mu(q,q)\tilde V_\alpha(q,0)\,D_\nu(q)C(q)^3$ up to the two phases of the internal photon, which cancel, with $\tilde V_\mu(q,q)=-\frac2a\sin(aq_\mu)$ odd and the other factors even; for the swordfish vertex (the seagull with one external photon leg, its other leg joined by the internal photon to a current vertex) the integrand is $(1+e^{-\mathrm iaq_\mu})(e^{\mathrm iaq_\mu}-1)D_\nu(q)C(q)^2=2\mathrm i\sin(aq_\mu)D_\nu(q)C(q)^2$, odd. ($\Gamma_8$ has no removed term, its $A^3$-tadpole containing no scalar line; the subtraction $t^0$ of that tadpole is $0$ anyway, as $V_3\propto\sin\frac{a(p+p')_\mu}2$.) Type (a) ($\Gamma_0$ with $l$ the inner line): $t^1_\gamma F_{\gamma,\mathrm{dbl}}=F_{\gamma,\mathrm{dbl}}(k_2;0)+P(k_1)\cdot\nabla F_{\gamma,\mathrm{dbl}}(k_2;0)$, and the order-$1$ term is odd in the inner loop momentum $k_2$ (as in the proof of \cref{lem:a05:forest}(a)), so the removed term's double integral is $N_\gamma\,I_{\bar\Gamma}$ with
\[
N_\gamma=\big\langle D_\nu(k)\,\hat e(k)\,C(k)^2\big\rangle_a,\qquad I_{\bar\Gamma}=\big\langle D_\nu(k)\,\hat e(k)\,C(k)^2\big\rangle_a,
\]
the inner exchange graph with its scalar line doubled at zero external momentum, and the outer exchange graph with the contracted subgraph as a unit vertex between its two scalar lines of equal momentum; here $\sum_\mu\tilde V_\mu(0,-k)\tilde V_\mu(-k,0)=\hat e(k)$. Since $D_\nu\hat e\le1$, both are $\le\langle C^2\rangle_a\le\frac{\pi^2}{64}$.

\emph{(3) Class B ($\Gamma_1$, $\Gamma_4$).} One photon chord $k$, one polarisation chord $q$. By \cref{lem:a05:PF} every term is $D_\nu(k)^2$ times a $\nu$-free function, which is the path factor times the subtracted bubble: the path factor is $\tilde V_\mu(0,-k)\tilde V_\nu(-k,0)C(k)$ for $\Gamma_1$ (with $C(k)^2$ when $l$ is the path line), bounded entrywise by $\hat e(k)C(k)\le1$, and a constant for $\Gamma_4$ (the seagull at zero scalar momenta and zero total photon momentum). \emph{Soft $k$:} use $R'F_{\mathrm{dbl}}=D_\nu(k)^2\,\mathrm{path}(k)\,(1-t^1_k)F_{b(,\mathrm{dbl})}(q;k)$ itself ($R'$ subtracts the bubble by definition; for $l$ on the path, $R'=R_{\mathrm{dbl}}$). By \cref{lem:a05:VPR} its $q$-integral is bounded by $C_{\mathrm{VP}}|k|^2|\mathrm{path}(k)|$, and $\langle\mathbf 1_{Q_s}D_\nu(k)^2|k|^2\rangle_a\le\frac{\pi^2}{64}$ by \eqref{eq:a05:soft}, uniformly in $\nu$. \emph{Hard $k$:} $R'F_{\mathrm{dbl}}=R_{\mathrm{dbl}}F_{\mathrm{dbl}}$ minus the removed term, which is present only when $l$ lies in the bubble. The $R_{\mathrm{dbl}}$ part is $\le(1+\pi^2)^2\langle\langle|R_{\mathrm{dbl}}F^{(1)}|\rangle\rangle_a\le(1+\pi^2)^2K$. The removed part is
\[
\big\langle\mathbf 1_{\mathrm{hard}}\,D_\nu(k)^2\,|\mathrm{path}(k)|\,\big|\Pi_{\mathrm{dbl}}(0)+P(k)\cdot\nabla_k\langle F_{b,\mathrm{dbl}}(\cdot;k)\rangle_a\big|_{k=0}\big|\big\rangle_{a,k}
\le(1+\pi^2)^2\,|\Pi_{\mathrm{dbl}}(0)|\,\big\langle D_1^2\big(1+Ca\sqrt{\hat e}\big)\big\rangle_a\le C\Big(1+a\big(1+\ln\tfrac1a\big)\Big)\le C,
\]
by \cref{lem:a05:VPR}, \cref{lem:a05:EB}(iv) and \eqref{eq:a05:soft} ($D_1^2\sqrt{\hat e}\le\chi^{-3/2}$). ($\Gamma_1$ vanishes identically on $\eps\Z^3$ by the transversality noted in \cref{lem:a05:40}(b); the bound covers it regardless.)

Summing (0)--(3) over the finitely many pairs $(\Gamma,l)$ gives $|Y(\nu,a)|\le C_U$, where $C_U$ collects $K$ and numerical constants, and depends on $a_0$ only.
\end{proof}

\begin{proof}[Proof of (v)]
By \cref{lem:a05:SC}, $y_{40}(s;\eps)=s^{-1}Y(\lmu/\sqrt s,\eps\sqrt s)$. For $\eps\le a_0/\sqrt s$ the second argument lies in $(0,a_0]$, and \cref{prop:a05:U} gives $|y_{40}(s;\eps)|\le C_U/s$ for every $\lmu>0$. So (v) holds with $C_{40}\defeq C_U$, which depends neither on $\lmu$ nor on $\eps$ nor on $s$. In scaled form, $s\,y_{40}$ is bounded by a function of $\lmu/\sqrt s$ and $\eps\sqrt s$ that stays bounded as $\lmu/\sqrt s\to0$: no $\ln(\sqrt s/\lmu)$ occurs.
\end{proof}

The mechanism of (iv) is the one of \cref{lem:a05:dec}: the only coefficient that grows with $s$, $|\lmu^2-2s|$, multiplies $B_s(p)-B_s(0)=O(|p|^2s^{-3/2})$, the Taylor step being taken on the scalar line; the linearly divergent constants cancel exactly in the decomposition (\cref{lem:a05:Kdec}) before any bound. In (v), $y_{40}(s;\eps)=s^{-1}Y(\lmu/\sqrt s,\eps\sqrt s)$ by scaling (\cref{lem:a05:SC}), and $Y$ is bounded uniformly in its first argument and in the rescaled spacing $a=\eps\sqrt s\le a_0$ (\cref{prop:a05:U}); the power-counting input enters at one mass configuration only, and the constants $C_{40}$, $a_0$ are those of \cref{lem:a05:H}.

\begin{theorem}[$k$-uniform lower bound]\label{thm:a05:A}
Let $e,\lambda\ge0$, $\lmu,m>0$, $M\ge m$ and $\eps_k\le a_0/M$. Then
\[
2u_k(M)\ \ge\ (M^2-m^2)-\frac{\pi^2}{16}e^2(M-m)-2\Big[C_{40}e^4+C_{21}(m^2)\lambda e^2+\frac{81\pi^4}{32}\lambda^2\Big]\ln\frac Mm-2C'_{21}\lambda e^2\eps_k(M-m).
\]
Consequently $s_-(M)\ge\frac12\Phi_A(M)$ with $\Phi_A(M)\defeq(M^2-m^2)-\frac{\pi^2}{16}e^2(M-m)-2[C_{40}e^4+C_{21}(m^2)\lambda e^2+\frac{81\pi^4}{32}\lambda^2]\ln\frac Mm\to\infty$ as $M\to\infty$; this is \NEGinf, and $s_-(M_2)>0$ at every $M_2$ with $\Phi_A(M_2)>0$.
\end{theorem}
\begin{proof}
Insert \cref{lem:a05:D123,lem:a05:D45} into \eqref{eq:a05:exact}: $y_{01}>0$ drops the $\lambda$ term; $\int_{m^2}^{M^2}|y_{20}|\,ds\le\int\frac{\pi^2}{32\sqrt s}\,ds=\frac{\pi^2}{16}(M-m)$; $\int|y_{02}|\le\frac{81\pi^4}{32}\ln\frac{M^2}{m^2}$; $\int|y_{21}|\le C_{21}(m^2)\ln\frac{M^2}{m^2}+2C'_{21}\eps(M-m)$, since $C_{21}(s)\le C_{21}(m^2)$ for $s\ge m^2$; $\int|y_{40}|\le C_{40}\ln\frac{M^2}{m^2}$, since $\eps\le a_0/M\le a_0/\sqrt s$ on $[m^2,M^2]$. Finally $\ln\frac{M^2}{m^2}=2\ln\frac Mm$, and the $\liminf$ drops the $\eps_k$ term.
\end{proof}
The constants of \cref{thm:a05:A} are crude, and $C_{40}$, $a_0$ are not explicit. The exact continuum limit, computed next, gives explicit and essentially sharp constants.

\paragraph{Continuum conventions.} $\int_k\defeq\int_{\R^3}\frac{d^3k}{(2\pi)^3}$. Complex field $\phi_c=(\phi_1+\mathrm i\phi_2)/\sqrt2$, so that $\frac12|D_A\phi|^2=|D_A\phi_c|^2$ and $\lambda|\phi|^4=4\lambda|\phi_c|^4$. Lines $G(q)=(q^2+M^2)^{-1}$, $D(k)=(k^2+\mu^2)^{-1}\delta_{\mu\nu}$ with $\mu\defeq\lmu$. Vertices (coefficients in the action): current $-e(p+p')_\mu$ ($p$ in, $p'$ out); seagull $2e^2\delta_{\mu\nu}$; quartic $16\lambda$; counterterm ${\dm}$. $\Sigma$ is the one-particle-irreducible self-energy, and King's rule is ${\dm}_{(\alpha\beta)}=\Sigma^{(\alpha\beta)}_{\rm int}(0)$. One-loop integrals: $I\defeq\int_kDG=\frac1{4\pi(M+\mu)}$; $J\defeq\int_kDq^2G^2=\frac{M+2\mu}{8\pi(M+\mu)^2}$; $\int G^2=\frac1{8\pi M}$; $\Bc_{ab}(q)=\int_kG_a(k)G_b(q+k)=\frac{\arctan(|q|/(a+b))}{4\pi|q|}$ (\cref{lem:a05:bub}(i)). Two-loop vacuum integrals with lines $q,k,q+k$ of masses $a,b,c$ and $\mathsf s\defeq a+b+c$: $V(a,b,c)\defeq\int\!\!\int G_a(q)G_b(k)G_c(q+k)$, of which only differences are used, $V(a,b,c)-V(a',b',c')=-\frac1{16\pi^2}\ln\frac{\mathsf s}{\mathsf s'}$ (absolutely convergent in momentum space, and equal to the position-space Frullani integral $\frac1{16\pi^2}\int_0^\infty(e^{-\mathsf sr}-e^{-\mathsf s'r})\frac{dr}r$ by Plancherel, since $\tilde G_a(x)=\frac{e^{-a|x|}}{4\pi|x|}$); and $V_2(a;b,c)\defeq\int\!\!\int G_a(q)^2G_b(k)G_c(q+k)=\frac1{32\pi^2a\,\mathsf s}$, because $\widetilde{(G_a^2)}(x)=\frac{e^{-a|x|}}{8\pi a}$. Write $V(\mathsf s)$ for $V$ with mass sum $\mathsf s$; only $V_1\defeq V(2\mu+M)$ and $V_2\defeq V(\mu+2M)$ occur. Two one-loop functions enter:
\begin{align}
\sigma(q)&\defeq\Sigma^{(20)}_{\rm ren}(q)/e^2=2q^2\Bc_{M\mu}(q)+(\mu^2-2M^2)\big(\Bc_{M\mu}(q)-\Bc_{M\mu}(0)\big)\ =\ 2u_{M,\mu}(q)+\tfrac{|q|}4,\label{eq:a05:sigma}\\
\Pi_{\mu\mu}(k)&=-e^2\big[k^2\Bc_{MM}(k)+4M^2\big(\Bc_{MM}(k)-\Bc_{MM}(0)\big)\big],\label{eq:a05:Pi}
\end{align}
the continuum renormalised order-$e^2$ self-energy (the continuum analogue after \cref{lem:a05:dec}, with $u$ as in \cref{thm:a05:Z}) and the trace of the transverse one-loop vacuum polarisation. For \eqref{eq:a05:Pi}: the current loop is $e^2\int(2l+k)^2GG'=e^2[4T-(k^2+4M^2)\Bc_{MM}]$ with $T\defeq\int G$, by $(2l+k)^2=2G^{-1}+2G'^{-1}-k^2-4M^2$; the seagull loop is $-6e^2T$, $k$-independent. On the lattice the full one-loop polarisation vanishes at $k=0$ exactly (the lattice Ward identity of King's gauge-invariant action), so the continuum limit is the subtracted form $\Pi^{\rm cur}(k)-\Pi^{\rm cur}(0)$; since $4\int l^2G^2=4T-4M^2\Bc_{MM}(0)$, $T$ cancels identically and \eqref{eq:a05:Pi} follows, with no regularisation. Check: $\Pi_{\mu\mu}(k)=-\frac{e^2k^2}{12\pi M}+O(k^4)$.

\begin{lemma}[the one-loop limits]\label{lem:a05:oneloop}
For all $M,m>0$, as $\eps\to0$: $d_{01}(M)-d_{01}(m)\to F_{01}(M)-F_{01}(m)$ and $d_{20}(M)-d_{20}(m)\to F_{20}(M)-F_{20}(m)$, where
\[
F_{01}(M)=\frac{4M}\pi,\qquad F_{20}(M)=-\frac{M^2}{4\pi(M+\mu)},\qquad D_{20}\defeq F_{20}(M)-F_{20}(m)=-\frac{(M-m)(Mm+M\mu+m\mu)}{4\pi(M+\mu)(m+\mu)}\in\Big[-\frac{M-m}{4\pi},0\Big]\ \ (M\ge m).
\]
\end{lemma}
\begin{proof}
$d_{01}(M)-d_{01}(m)=16(M^2-m^2)\langle C_MC_m\rangle$ and $d_{20}(M)-d_{20}(m)=-(M^2-m^2)\langle\hat q^2DC_MC_m\rangle$. On $\BZ_\eps$, $\hat q^2\ge\frac4{\pi^2}q^2$ and $\hat q^2D\le1$, so both integrands are dominated by $(\frac{\pi^2}4)^2\big((q^2+M^2)(q^2+m^2)\big)^{-1}$, integrable on $\R^3$, and converge pointwise. Dominated convergence: $16(M^2-m^2)\int G_MG_m=\frac{16(M-m)}{4\pi}$, and $(M^2-m^2)\int k^2DG_MG_m=\frac{(M-m)(Mm+M\mu+m\mu)}{4\pi(M+\mu)(m+\mu)}$ via $\frac{k^2}{(k^2+a^2)(k^2+b^2)}=\frac1{b^2-a^2}\big[\frac{b^2}{k^2+b^2}-\frac{a^2}{k^2+a^2}\big]$, which is $F_{20}(M)-F_{20}(m)$. The bracket follows from $D_{20}=-\frac{M-m}{4\pi}\big(1-\frac{\mu^2}{(M+\mu)(m+\mu)}\big)$.
\end{proof}

\begin{lemma}[the sunset]\label{lem:a05:02}
For all $M,m>0$, $d_{02}(M)-d_{02}(m)\to D_{02}\defeq F_{02}(M)-F_{02}(m)$ with $F_{02}(M)=-\frac8{\pi^2}\ln M$; the $M^2$-derivative is $y^\infty_{02}(s)=-\frac4{\pi^2s}$.
\end{lemma}
\begin{proof}
For each $s>0$, $y_{02}(s;\eps)\to-384\int\!\!\int G(q)^2G(r)G(q+r)=-384\,V_2(M;M,M)=-\frac{384}{96\pi^2M^2}=-\frac4{\pi^2s}$: the lattice integrand, extended by zero outside $\BZ_\eps^2$, is dominated by the integrable function $(\frac{\pi^2}4)^4\sum_P\tilde G(q)^2\tilde G(r)\tilde G(q+r-P)$ of the proof of \cref{lem:a05:D123}(iii) and converges pointwise (for fixed $(q,r)$ eventually $q+r\in\BZ_\eps$, so $P=0$). By \cref{lem:a05:D123}(iii), $|y_{02}(s;\eps)|\le\frac{81\pi^4}{32s}$ uniformly in $\eps$, so dominated convergence in $s$ on $[m^2,M^2]$ (or $[M^2,m^2]$) in \eqref{eq:a05:exact} gives $d_{02}(M)-d_{02}(m)\to\int_{m^2}^{M^2}(-\frac4{\pi^2s})\,ds=-\frac8{\pi^2}\ln\frac Mm$.
\end{proof}

\begin{lemma}[the mixed order]\label{lem:a05:21}
For all $M,m>0$, $d_{21}(M)-d_{21}(m)\to D_{21}\defeq F_{21}(M)-F_{21}(m)$ with
\[
F_{21}(M)=\frac2{\pi^2}\ln(2M+\mu)+\frac{\mu}{2\pi^2(M+\mu)},\qquad
y^\infty_{21}(s)=\frac{dF_{21}}{dM^2}=\frac{8M^2+14M\mu+7\mu^2}{4\pi^2M(M+\mu)^2(2M+\mu)}>0 .
\]
In particular $D_{21}\ge0$ for $M\ge m$.
\end{lemma}
\begin{proof}
\emph{Limit.} $d_{21}=-8X_\eps(M;\mu)$, and \cref{thm:a05:Z} with $t=\eps$ at the masses $(M,\mu)$ and $(m,\mu)$ gives $d_{21}(M)-d_{21}(m)\to-8\big[\kappa_{\rm c}(M,\mu)-\kappa_{\rm c}(m,\mu)\big]$, the $\ln\frac1\eps$ terms and $\kappa_{\rm lat}$ cancelling. No input beyond \cref{thm:a05:Z} is used.
\emph{Value.} By \eqref{eq:a05:sigma}, $u_{M,\mu}=\frac12\sigma-\frac{|q|}8$, so $\kappa_{\rm c}(M,\mu)=2\int_qG^2\sigma-\frac1{2\pi^3}\int\frac{|q|/8}{(q^2+M^2)^2}d^3q-\frac{\ln M+\frac12}{4\pi^2}$ in the sense of $M^2$-derivatives, each of which is absolutely convergent; the derivative of the last two terms is $\frac1{8\pi^2M^2}-\frac1{8\pi^2M^2}=0$. So, up to an $M$-independent constant, $-8\kappa_{\rm c}=-16\int G^2\sigma$. Partial fractions in $q^2$, with $q^2G^2=G-M^2G^2$:
$\int G^2\sigma=2\big[V(M,\mu,M)-M^2V_2(M;\mu,M)\big]+(\mu^2-2M^2)\big[V_2(M;\mu,M)-\frac{\Bc_{M\mu}(0)}{8\pi M}\big]$, with $V(M,\mu,M)=V_2$ and $V_2(M;\mu,M)=\frac1{32\pi^2M(2M+\mu)}$. Hence
\[
-16\int G^2\sigma=-32V_2+\frac1{2\pi^2}\Big[\frac{4M^2-\mu^2}{M(2M+\mu)}-\frac{2M^2-\mu^2}{M(M+\mu)}\Big]=\frac2{\pi^2}\ln(2M+\mu)+\frac{\mu}{2\pi^2(M+\mu)}+\text{const},
\]
since the bracket equals $\frac{\mu}{M+\mu}$. Each identity holds between the $M^2$-derivatives, which are absolutely convergent (Fubini). Differentiating $F_{21}$ gives $y^\infty_{21}$; its numerator is positive.
\end{proof}

\begin{lemma}[the order $e^4$]\label{lem:a05:40}
For all $M,m>0$, $d_{40}(M)-d_{40}(m)\to D_{40}\defeq F_{40}(M)-F_{40}(m)$, where
\begin{gather*}
F_{40}(M)=-\frac1{4\pi^2}\ln M+\Phi_{40}\Big(\frac\mu M\Big),\\
\Phi_{40}(x)=\frac{-7+4x+8x^2-3x^3-3x^4+x^5}{32\pi^2(1+x)(1-x^2)^2}+\frac{(5-2x^2)\ln(2+x)-(9-10x^2+4x^4)\ln(1+2x)}{16\pi^2(1-x^2)^2},
\end{gather*}
continuous on $[0,\infty)$ (also at $x=1$), with $\Phi_{40}(0)=\frac{10\ln2-7}{32\pi^2}$ and $\Phi_{40}(x)=-\frac1{4\pi^2}\ln x+O(1)$ as $x\to\infty$. The $M^2$-derivative is $y^\infty_{40}=-\frac1{8\pi^2M^2}-\frac{\mu}{2M^3}\Phi_{40}'(\frac\mu M)$.
\end{lemma}
\begin{proof}
\emph{(a) Existence of the limit.} Fix $s>0$. The integrand of $y_{40}(s;\eps)$ is a sum of two-loop lattice Feynman integrands of the class treated in \cite{Rei88,Rei88b}: standard lattice propagators with positive masses, trigonometric-polynomial numerators from the vertices of King's action (the current vertex $V_1$, the seagull $V_2$, and the three- and four-photon vertices $V_3$, $V_4$ of the lattice hopping term), and the insertion of $d_{20}$ on the exchange line, which is exactly the zero-momentum Taylor subtraction of the order-$e^2$ two-point subgraph. The remaining divergent subgraphs are the logarithmically divergent one-loop vertex subgraphs (triangle and swordfish); their zero-momentum Taylor terms vanish identically, each being odd under the reflection of its loop momentum by the oddness of the current vertex, with the two phases of the internal photon cancelling, so these subtractions may be inserted without changing the integrand. The integrand is then the subtracted integrand of the forest formula of \cite{Rei88b}, whose lattice degree of divergence, after the derivative in $s$, is negative on every Zimmermann subspace (the degree count of \cite{Rei88} in three dimensions, to which it extends as remarked there), the loop momenta being routed along the photon chords so that the set of subtracted subgraphs is natural in the sense of \cite{Rei88}. The power-counting theorem of \cite{Rei88} gives the convergence of $y_{40}(s;\eps)$ to the integral over $\R^6$ of the pointwise limit of the subtracted integrand. In that limit the lattice-only structures vanish: $V_3$ and $V_4$ carry $\eps$ and $\eps^2$ relative to $V_1$ (their photon tadpoles are $O(\eps^{-1})$, so the dressed vertices are still $O(\eps)$), the cosine factors of the seagull tend to $1$, and the phases tend to $1$. So $y^\infty_{40}(s)$ is the $s$-derivative of King's continuum graph classes with the subtractions above, and by \cref{lem:a05:D45}(v), for $\eps\le a_0/\max(M,m)$, $|y_{40}(s;\eps)|\le C_{40}/s$ on the interval between $m^2$ and $M^2$, so dominated convergence in \eqref{eq:a05:exact} gives $d_{40}(M)-d_{40}(m)\to\int_{m^2}^{M^2}y^\infty_{40}(s)\,ds$.

\emph{(b) The classes.} The twelve two-loop one-particle-irreducible two-point topologies of order $e^4$ fall into the continuum classes of the following table (values per $e^4$, up to $M$-independent constants; $k_3=k_1+k_2$, $G_i=G(k_i)$, $D_i=D(k_i)$):
\begin{center}
\begin{tabular}{@{}lll@{}}
topology & class & value\\\hline
rainbow, and seagull tadpole on the inner line, $d_{20}$ inserted & $\mathrm R$ & $\int D(q)q^2G(q)^2\sigma(q)$\\
crossed & $\mathrm X$ & $V_1-M^2I^2$\\
seagull at an end, both photons to the currents (two mirror graphs) & $\mathrm{SV}_2$ & $-2[(M^2-\mu^2)I^2+3V_1+(\mu^2-4M^2)W_2]$\\
two-photon seagull sunset (melon) & $\mathrm{Me}$ & $6V_1$\\
seagull tadpole with the one-loop polarisation inserted & $\mathrm{SS}$ & $V_2-\frac1{32\pi^2}$\\
exchange with the polarisation inserted & & $0$\\
seagull in the middle & & $0$\\
$V_3$-tadpole vertex & & $0$\\
$V_4$ double tadpole & & $M$-independent
\end{tabular}
\end{center}
Here $W_2\defeq\int\!\!\int D_1D_2G_2G_3=\frac{V_1-V_2}{M^2-\mu^2}$, by $DG=\frac{D-G}{M^2-\mu^2}$. The last four rows: the polarisation insertions on the exchange line vanish identically on $\eps\Z^3$, because the lattice current vertex at zero scalar momentum is $-\mathrm i\hat k'_\mu(k)$ and $\hat k'_\mu\Pi_{\mu\nu}(k)=0$ (the lattice Ward identity; in the continuum $k_\mu\Pi_{\mu\nu}=0$). The seagull in the middle factorises, at zero external momentum, into $(\int D(k)k_\mu G(k))^2$, a product of two odd integrals, so $0$ in the continuum, and is $O(\eps)$ on the lattice; the $V_3$ class is $O(\eps)$ as said in (a); the $V_4$ double tadpole has no scalar line, so its $s$-derivative vanishes. The inner seagull tadpole of the second graph of the first row is momentum dependent on the lattice only at $O(\eps q^2)$; its sum with the rainbow and the counterterm is the lattice renormalised self-energy, whose continuum limit is $\sigma$.

\emph{(c) The class values} (continuum, zero external momentum; signs and combinatorics from the Feynman rules above).
$\mathrm X=e^4\int\!\!\int D_1D_2G_1G_2G_3\,[k_1\cdot(k_1+2k_2)][k_2\cdot(2k_1+k_2)]$. The exact identities $k_1\cdot(k_1+2k_2)=k_3^2-k_2^2=G_3^{-1}-G_2^{-1}$ and $k_2\cdot(2k_1+k_2)=G_3^{-1}-G_1^{-1}$ turn the integrand into $D_1D_2[G_1G_2G_3^{-1}-G_1-G_2+G_3]$ with $G_3^{-1}=k_1^2+k_2^2+2k_1\cdot k_2+M^2$; the products factorise, the $k_1\cdot k_2$ term is odd, and with $\int Dk^2G=T_\mu-M^2I$ (pointwise $DGk^2=D-M^2DG$, $T_\mu\defeq\int D$) the $T_\mu$ terms cancel identically: $\mathrm X/e^4=2I(T_\mu-M^2I)+M^2I^2-2IT_\mu+V_1=V_1-M^2I^2$.
$\mathrm{SV}_2$ (both mirror graphs): $2\mathrm{SV}=-4e^4\int\!\!\int D_1D_2G_2G_3\,(k_1+2k_2)\cdot k_2$; with $(k_1+2k_2)\cdot k_2=\frac12(G_3^{-1}+3G_2^{-1}-D_1^{-1}-4M^2+\mu^2)$ the four terms give $T_\mu I$, $3V_1$, $-IT_M$ ($T_M\defeq\int G$) and $(\mu^2-4M^2)W_2$, and $T_\mu-T_M=(M^2-\mu^2)I$.
$\mathrm{Me}=\frac12(2e^2)^2\delta_{\mu\nu}\delta_{\mu\nu}\int\!\!\int DDG=6e^4V_1$; its $M^2$-derivative is $-6e^4\int\!\!\int D_1D_2G_3^2=-\frac{6e^4}{32\pi^2M(M+2\mu)}$.
$\mathrm{SS}=-e^2\int D(k)^2\Pi_{\mu\mu}(k)=e^4\int D^2[k^2\Bc_{MM}+4M^2(\Bc_{MM}-\Bc_{MM}(0))]$; with $k^2=D^{-1}-\mu^2$ and $\int D^2=\frac1{8\pi\mu}$: $\mathrm{SS}/e^4=V_2-\mu^2V_2(\mu;M,M)+4M^2\big[V_2(\mu;M,M)-\frac1{8\pi M}\cdot\frac1{8\pi\mu}\big]=V_2+\frac{2M-\mu}{32\pi^2\mu}-\frac M{16\pi^2\mu}=V_2-\frac1{32\pi^2}$; the $1/\mu$ terms cancel.
$\mathrm R=e^4\int Dq^2G^2\sigma=e^4\big[2\int Dq^4G^2\Bc+(\mu^2-2M^2)\int Dq^2G^2(\Bc-\Bc(0))\big]$, $\Bc=\Bc_{M\mu}$. Partial fractions: $\frac{x^2}{(x+\mu^2)(x+M^2)^2}=\frac a{x+\mu^2}+\frac{1-a}{x+M^2}+\frac c{(x+M^2)^2}$ with $a=\frac{\mu^4}{(M^2-\mu^2)^2}$, $c=-\frac{M^4}{M^2-\mu^2}$, and $\frac x{(x+\mu^2)(x+M^2)^2}=\frac A{x+\mu^2}-\frac A{x+M^2}+\frac{C'}{(x+M^2)^2}$ with $A=-\frac{\mu^2}{(M^2-\mu^2)^2}$, $C'=\frac{M^2}{M^2-\mu^2}$. So
\[
\mathrm R/e^4=2\big[aV_1+(1-a)V_2+c\,V_{2M}\big]+(\mu^2-2M^2)\big[A(V_1-V_2)+C'V_{2M}-JI\big],\qquad V_{2M}\defeq V_2(M;\mu,M)=\frac1{32\pi^2M(2M+\mu)}.
\]
No linear divergence is cancelled between the pieces: $\sigma$ is convergent.

\emph{(d) Why these manipulations give the subtracted limit of (a).} $\mathrm R$, $\mathrm{SS}$ and $\mathrm{Me}$ use only the convergent functions $\sigma$ and $\Pi$, partial fractions and $V$-differences; their $M^2$-derivatives are absolutely convergent integrals, and the identities hold between them. $\mathrm X$ and $\mathrm{SV}_2$ contain a logarithmically divergent vertex subgraph. For these use the product cutoff $|k_1|,|k_2|<\Lambda$ in the photon-chord routing. Every subtracted subgraph's loop momentum is then a cutoff variable, and its zero-momentum Taylor term is odd in that variable, so it integrates to zero on the symmetric ball; hence the product-cutoff integral of the unsubtracted class equals that of its subtracted integrand, which converges to the subtracted value by dominated convergence. The reduction of $\mathrm X$ uses no shift of integration variables and is exact at every $\Lambda$. The term $\int\!\!\int D_2G_2G(k_1+k_2)$ of $\mathrm{SV}_2$ needs one shift of the linearly divergent $k_1$-subintegral; the surface term $E_\Lambda=\int_{|k_2|<\Lambda}D_2G_2\,\delta_\Lambda(k_2)$, $\delta_\Lambda(k_2)\defeq\int_{|k_1|<\Lambda}[G(k_1)-G(k_1+k_2)]$, has $dE_\Lambda/dM^2\to0$, since $|\delta_\Lambda(k)|\le C|k|^2/\Lambda$ for $|k|\le\Lambda/2$ and $\le C\Lambda$ otherwise, and the $M^2$-derivative brings one more factor $G$. Hence the $M^2$-derivatives of $\mathrm X$ and $\mathrm{SV}_2$ are as displayed.

\emph{(e) The sum.} $F_{40}=\mathrm X+\mathrm{Me}+\mathrm{SV}_2+\mathrm{SS}+\mathrm R$. The coefficients of the logarithmically divergent $V$'s are $1+6-6+1+2=4$, so $F_{40}$ carries $\frac4{16\pi^2}\ln\frac1\eps=\frac1{4\pi^2}\ln\frac1\eps$, the universal two-loop coefficient of order $e^4$; the $M$-dependence is finite. Substituting the closed forms of $V_1-V_2$, $W_2$, $I$, $J$, $V_{2M}$ and $V_2(\mu;M,M)$ and simplifying gives $F_{40}(M,\mu)=-\frac1{4\pi^2}\ln M+\Phi_{40}(\mu/M)$ as an exact identity between elementary functions; $\Phi_{40}$ is continuous at $x=1$ (the apparent poles of the two terms cancel), $\Phi_{40}(0)=\frac{10\ln2-7}{32\pi^2}$, and $\Phi_{40}(x)=-\frac1{4\pi^2}\ln x+O(1)$ as $x\to\infty$ (decoupling). At $\mu\to0$, $F_{40}(M)-F_{40}(m)\to-\frac1{4\pi^2}\ln\frac Mm$ exactly: the massless-photon limit is infrared finite. The derivative formula is $\frac{d}{dM^2}F_{40}=\frac1{2M}\frac d{dM}F_{40}$.
\end{proof}

\begin{theorem}[the exact continuum limit of the bare coefficient]\label{thm:a05:E}
For all $M,m>0$ the limit $u_\infty(M)\defeq\lim_ku_k(M)$ exists, and
\[
2u_\infty(M)=(M^2-m^2)+\lambda D_{01}+e^2D_{20}+e^4D_{40}+\lambda e^2D_{21}+\lambda^2D_{02},\qquad D_{\alpha\beta}\defeq F_{\alpha\beta}(M)-F_{\alpha\beta}(m),
\]
with $F_{01}$, $F_{20}$ of \cref{lem:a05:oneloop}, $F_{02}$ of \cref{lem:a05:02}, $F_{21}$ of \cref{lem:a05:21} and $F_{40}$ of \cref{lem:a05:40}. The $M^2$-derivatives are
\[
y_{01}=\frac2{\pi M},\quad y_{20}=-\frac{M+2\mu}{8\pi(M+\mu)^2},\quad y_{02}=-\frac4{\pi^2M^2},\quad y_{21}=\frac{8M^2+14M\mu+7\mu^2}{4\pi^2M(M+\mu)^2(2M+\mu)}>0,\quad y_{40}=-\frac1{8\pi^2M^2}-\frac\mu{2M^3}\Phi_{40}'\Big(\frac\mu M\Big),
\]
and for large $M$, $2u_\infty(M)=M^2+\big(\frac{4\lambda}\pi-\frac{e^2}{4\pi}\big)M-\big(\frac{e^4}{4\pi^2}-\frac{2\lambda e^2}{\pi^2}+\frac{8\lambda^2}{\pi^2}\big)\ln M+O(1)$; the coefficient of $\ln M$ is minus the coefficient of $\ln\frac1\eps$ in the two-loop counterterm.
\end{theorem}
\begin{proof}
\Cref{lem:a05:oneloop,lem:a05:02,lem:a05:21,lem:a05:40} in \eqref{eq:a05:exact} at $\eps=\eps_k$.
\end{proof}

\begin{lemma}[monotonicity of $\Phi_{40}$]\label{lem:a05:S}
$\Phi_{40}'(x)<0$ for every $x\ge0$. Consequently, for $M\ge m$, $D_{40}\ge-\frac1{4\pi^2}\ln\frac Mm$; and $y_{40}>-\frac1{8\pi^2M^2}$ for all $M,\mu>0$.
\end{lemma}
\begin{proof}
(i) \emph{Derivative.} Differentiating the closed form,
\[
\Phi_{40}'(x)=\frac{N(x)}{32\pi^2(1-x^2)^3(1+x)^2(2+x)(1+2x)},\qquad N(x)=(x^2-1)Q(x)+P(x)\ln\frac{1+2x}{2+x},
\]
with $Q(x)=16x^7+56x^6+24x^5-100x^4-83x^3+82x^2+109x+40$ and $P(x)=8x(x-2)(x+1)^2(x+2)^2(2x+1)$. The denominator has the sign of $1-x$, so $\Phi_{40}'<0$ iff $\operatorname{sign}N(x)=\operatorname{sign}(x-1)$.
(ii) \emph{Substitution} $x=\frac{1+t}{1-t}$, $t\in[-1,1)$. Then $\ln\frac{1+2x}{2+x}=\ln\frac{3+t}{3-t}=2\operatorname{artanh}\frac t3$, and $\tilde N(t)\defeq N(x(t))(1-t)^9=A(t)+2B(t)\operatorname{artanh}\frac t3$ with
\[
A(t)=576t-2496t^2+13264t^3-3056t^4+224t^5-544t^6+272t^7-48t^8,\qquad B(t)=32(t-3)^2(t-1)^2(t+1)(t+3)(3t-1).
\]
Since $(1-t)^9>0$ and $\operatorname{sign}(x-1)=\operatorname{sign}t$, the claim is $\tilde N(t)/t^3>0$ on $[-1,1]$.
(iii) \emph{Exact series.} With $2\operatorname{artanh}\frac t3=\sum_{j\ge0}\frac{2t^{2j+1}}{3^{2j+1}(2j+1)}$ truncated at $j=30$, the coefficients of $t^0,t^1,t^2$ of $\tilde N$ vanish exactly, and $\tilde N(t)/t^3=S(t)+E(t)$, where $S$ is a polynomial with constant term $c_0=\frac{34160}3$ and $\sum_{k\ge1}|c_k|\le9103.93$ (exact rational arithmetic), and the tail obeys $|E(t)|\le2\max_{|t|\le1}|B|\cdot\frac{(1/3)^{63}}{63\cdot(8/9)}\le4.9\cdot10^{-28}$, with $\max|B|\le15808$ (the sum of the absolute values of its coefficients). Hence $\tilde N(t)/t^3\ge c_0-9103.93-4.9\cdot10^{-28}>2282>0$ on $[-1,1]$.
(iv) \emph{Uniform formula.} With $x^2-1=\frac{4t}{(1-t)^2}$, $(x^2-1)^3=\frac{64t^3}{(1-t)^6}$, so for $t\ne0$
\[
\Phi_{40}'(x(t))=-\frac{\tilde N(t)/t^3}{2048\pi^2(1-t)^3(1+x)^2(2+x)(1+2x)},
\]
and at $t=0$ ($x=1$) by continuity $\Phi_{40}'(1)=-\frac{c_0}{2048\pi^2\cdot36}\approx-0.0157$. The denominator is positive on $[-1,1)$ and $\tilde N/t^3>2282$.
\emph{Consequences.} For $M\ge m$, $\mu/M\le\mu/m$, so $\Phi_{40}(\mu/M)\ge\Phi_{40}(\mu/m)$, which is $D_{40}\ge-\frac1{4\pi^2}\ln\frac Mm$; and $y_{40}>-\frac1{8\pi^2M^2}$ by the derivative formula of \cref{lem:a05:40}.
\end{proof}

\begin{theorem}[the explicit lower bound]\label{thm:a05:B}
For every $M\ge m$ and every $\lmu>0$,
\[
2u_\infty(M)\ \ge\ (M-m)\Big(M+m+\frac{4\lambda}\pi-\frac{e^2}{4\pi}\Big)-\Big(\frac{e^4}{4\pi^2}+\frac{8\lambda^2}{\pi^2}\Big)\ln\frac Mm .
\]
With $M_2\defeq\max\{m,\ \frac{e^2}{4\pi}+(\frac{e^4}{4\pi^2}+\frac{8\lambda^2}{\pi^2})/m\}+1$: $s_-(M_2)=u_\infty(M_2)\ge\frac12$, and $s_-(M)\ge\frac12(M-m)(M-M_2+1)\to\infty$ for $M\ge M_2$.
\end{theorem}
\begin{proof}
Term by term in \cref{thm:a05:E}, for $M\ge m$: $\lambda D_{01}=\frac{4\lambda}\pi(M-m)$; $e^2D_{20}\ge-\frac{e^2}{4\pi}(M-m)$ (\cref{lem:a05:oneloop}); $e^4D_{40}\ge-\frac{e^4}{4\pi^2}\ln\frac Mm$ (\cref{lem:a05:S}); $\lambda e^2D_{21}\ge0$ (\cref{lem:a05:21}); $\lambda^2D_{02}=-\frac{8\lambda^2}{\pi^2}\ln\frac Mm$. This is the bound. With $\ln\frac Mm\le\frac{M-m}m$, $2u_\infty(M)\ge(M-m)\big[M+m+\frac{4\lambda}\pi-\frac{e^2}{4\pi}-(\frac{e^4}{4\pi^2}+\frac{8\lambda^2}{\pi^2})/m\big]\ge(M-m)(M-M_2+1)$ for $M\ge M_2$; at $M=M_2$ this is $\ge M_2-m\ge1$. Finally $s_-(M)=u_\infty(M)$, since the limit exists.
\end{proof}
The ``$+1$'' and the ``$\frac12$'' are in the mass unit of $m$, $e^2$, $\lambda$; replacing $+1$ by $+\delta$ gives $u_\infty(M_2(\delta))\ge\delta^2/2$ by the same proof.

\begin{proof}[Proof of \cref{prop:k02:NEG}]
\Cref{thm:a05:E} gives the existence of $u_\infty(M)$ for every $M>0$, \cref{thm:a05:B} the displayed inequality, the divergence $s_-(M)=u_\infty(M)\to\infty$ and the value $s_-(M_2)\ge\frac12$ at the stated $M_2$. The statement holds at every point of King's range, $e,\lambda\ge0$ and $m,\lmu>0$, with no exceptional set; \cref{thm:a05:A} gives \NEGinf\ independently of the closed forms. In particular $m$ is never a global maximiser of $M\mapsto\lim_k[\alpha_k(M)-\alpha_k(m)]$. At strong coupling $u_\infty(M)$ can be negative for $M$ up to $O(\frac{e^2}{4\pi}+\frac{e^4}{4\pi^2m})$ and is positive beyond.
\end{proof}

\begin{remark}\label{rem:a05:T2}
\Cref{thm:a05:E} holds for every $M>0$, on either side of $m$. Hence for every auxiliary mass $M_1>0$ the quantity $\sigma_1(M_1)=\liminf_k\frac12(\alpha_k(m)-\alpha_k(M_1))$ of \cref{hyp:s03:T2} is the limit $-u_\infty(M_1)$, finite and explicit; its one-loop part is the formula of \cref{lem:s14:CW}(b)(ii), and its two-loop part is $-\frac12[e^4D_{40}+\lambda e^2D_{21}+\lambda^2D_{02}]$ with the closed forms above. The threshold $s_*=\frac12\sup_M(\alpha_\infty(m)-\alpha_\infty(M))$ of \RW\ is thereby explicit as well.
\end{remark}
